%%%%%%%% ICML 2026 SUBMISSION, built on the official example_paper.tex %%%%%%%%%%%%%%%%%%
%
% Build:   ./make_figures.sh && latexmk -pdf main.tex   (or ./make_figures.sh --pdf)
%          The figures are not in git: make_figures.sh rebuilds figures/*.pdf from
%          scripts/ and data/ after every pull (see README.md).
% Status:  the main-text sections are an empty skeleton on purpose; the drafted material
%          lives in the appendices, to be moved into the main text section by section.

\documentclass{article}

% Recommended, but optional, packages for figures and better typesetting:
\usepackage{microtype}
\usepackage{graphicx}
\usepackage{subcaption}
\usepackage{booktabs} % for professional tables

% hyperref makes hyperlinks in the resulting PDF.
% If your build breaks (sometimes temporarily if a hyperlink spans a page)
% please comment out the following usepackage line and replace
% \usepackage{icml2026} with \usepackage[nohyperref]{icml2026} above.
\usepackage{hyperref}

% Attempt to make hyperref and algorithmic work together better:
\newcommand{\theHalgorithm}{\arabic{algorithm}}

% Use the following line for the initial blind version submitted for review:
\usepackage{icml2026}

% For preprint, use
% \usepackage[preprint]{icml2026}

% If accepted, instead use the following line for the camera-ready submission:
% \usepackage[accepted]{icml2026}

\usepackage{amsmath}
\usepackage{amssymb}
\usepackage{mathtools}
\usepackage{amsthm}
\usepackage{xcolor}
\usepackage{multirow}

% if you use cleveref..
\usepackage[capitalize,noabbrev]{cleveref}

%%%%%%%%%%%%%%%%%%%%%%%%%%%%%%%%
% THEOREMS
%%%%%%%%%%%%%%%%%%%%%%%%%%%%%%%%
\theoremstyle{plain}
\newtheorem{theorem}{Theorem}[section]
\newtheorem{proposition}[theorem]{Proposition}
\newtheorem{lemma}[theorem]{Lemma}
\newtheorem{corollary}[theorem]{Corollary}
\theoremstyle{definition}
\newtheorem{definition}[theorem]{Definition}
\newtheorem{assumption}[theorem]{Assumption}
\theoremstyle{remark}
\newtheorem{remark}[theorem]{Remark}

% Todonotes is useful during development; simply uncomment the next line
%    and comment out the line below the next line to turn off comments
%\usepackage[disable,textsize=tiny]{todonotes}
\usepackage[textsize=tiny]{todonotes}

% ---------------------------------------------------------------------------------------
% notation
% ---------------------------------------------------------------------------------------
\newcommand{\Mc}{\mathcal{M}}               % chirp mass
\newcommand{\Amp}{\mathcal{A}}              % strain amplitude
\newcommand{\Tobs}{T_{\rm obs}}             % observation time
\newcommand{\fw}{f_{w}}                     % centre frequency of a window
\newcommand{\fdot}{\dot f}
\newcommand{\snr}{\rho}
\newcommand{\Msun}{M_\odot}
\newcommand{\R}{\mathbb{R}}
\newcommand{\Sph}{\mathbb{S}}
\newcommand{\E}{\mathbb{E}}
\newcommand{\Var}{\operatorname{Var}}
\newcommand{\MAD}{\operatorname{MAD}}
\newcommand{\Exp}{\operatorname{Exp}}
\newcommand{\Log}{\operatorname{Log}}
\newcommand{\Unif}{\mathcal{U}}
\newcommand{\Normal}{\mathcal{N}}
\newcommand{\CN}{\mathcal{CN}}
\newcommand{\dd}{\mathrm{d}}
\newcommand{\ii}{\mathrm{i}}
\newcommand{\T}{^{\top}}
\newcommand{\inner}[2]{\langle #1 \,|\, #2 \rangle}
\newcommand{\bins}{\,\text{bins}}
\newcommand{\manifold}{\mathcal{X}}         % one source's flow manifold
\newcommand{\vnet}{\hat v_{\eta}}           % the network's velocity

% ---------------------------------------------------------------------------------------
% drafting aids: remove before submission
% ---------------------------------------------------------------------------------------
\definecolor{tbdcolor}{HTML}{C4351F}
% a value that B1 (or another pending run) will fill in
\newcommand{\tbd}[1][]{\textcolor{tbdcolor}{\textbf{[TBD#1]}}}
% a figure slot that a pending run will fill
\newcommand{\placeholderfig}[2][0.32\textwidth]{%
  \fbox{\parbox[c][#1][c]{0.95\linewidth}{\centering\small\color{tbdcolor}
  \textbf{placeholder}\\[2pt] #2}}}
% notes to the authors, rendered in the margin while drafting
\newcommand{\draftnote}[1]{\todo[color=yellow!25, size=\tiny]{#1}}


% The \icmltitle you define below is probably too long as a header.
% Therefore, a short form for the running title is supplied here:
\icmltitlerunning{Flow Matching Posteriors for LISA Galactic Binaries}

\begin{document}

\twocolumn[
  % working title, to be settled
  \icmltitle{Amortised Posteriors for LISA Galactic Binaries\\
    with Flow Matching on a Product Manifold}

  % It is OKAY to include author information, even for blind submissions: the
  % style file will automatically remove it for you unless you've provided
  % the [accepted] option to the icml2026 package.

  % TODO: author list and affiliations
  \begin{icmlauthorlist}
    \icmlauthor{Giorgio Mentasti}{apc}
    \icmlauthor{Firstname2 Lastname2}{yyy}
  \end{icmlauthorlist}

  \icmlaffiliation{apc}{Universit\'e Paris Cit\'e, CNRS, Astroparticule et Cosmologie, F-75013 Paris, France}
  \icmlaffiliation{yyy}{Department of XXX, University of YYY, Location, Country}

  % TODO: corresponding author and email
  \icmlcorrespondingauthor{Firstname1 Lastname1}{email@domain}

  % You may provide any keywords that you find helpful for describing your
  % paper; these are used to populate the "keywords" metadata in the PDF but
  % will not be shown in the document
  \icmlkeywords{flow matching, simulation-based inference, gravitational waves, LISA,
    Riemannian manifolds, amortised inference}

  \vskip 0.3in
]

% this must go after the closing bracket ] following \twocolumn[ ...

% This command actually creates the footnote in the first column listing the
% affiliations and the copyright notice. The command takes one argument, which
% is text to display at the start of the footnote. The \icmlEqualContribution
% command is standard text for equal contribution. Remove it (just {}) if you
% do not need this facility.

% Use ONE of the following lines. DO NOT remove the command.
% If you have no special notice, KEEP empty braces:
\printAffiliationsAndNotice{}  % no special notice (required even if empty)
% Or, if applicable, use the standard equal contribution text:
% \printAffiliationsAndNotice{\icmlEqualContribution}

\begin{abstract}
  % TODO: one paragraph, 4--6 sentences, once the main text is written.
  \tbd[ abstract]
\end{abstract}

% =======================================================================================
\section{Introduction}
\label{sec:introduction}
% =======================================================================================
% Planned paragraphs (empty on purpose; to be written from the appendices):
%  - diffusion models and flow matching as conditional samplers; amortised
%    simulation-based inference (NPE, FMPE, score-based SBI)
%  - LISA and its Galactic binaries: tens of millions of sources, ~10^4 resolvable,
%    overlapping in frequency; the population matters for astrophysics
%  - the global fit: trans-dimensional MCMC over all sources at once, its cost, and why
%    an amortised, local (per-window) posterior estimator is attractive
%  - what is hard for a generative model here: unordered sets of sources (label
%    switching), angles and spheres, exact degeneracies, posteriors 10^3 times narrower
%    than the prior at high SNR
%  - contributions: (i) flow matching on a product manifold for an unordered set of
%    binaries; (ii) why the flow-matching loss is blind to narrow posteriors and a time
%    warp that fixes it (app:sensitivity); (iii) a training recipe (warp, short auxiliary phase,
%    lr cooldown) and an ablation ladder on XS (app:xs); (iv) the whole LISA band (B,
%    app:b) and a comparison with MCMC (app:mcmc)

% =======================================================================================
\section{Galactic Binaries in LISA Data}
\label{sec:problem}
% =======================================================================================
% The physics, and the simulated dataset. Drafted material: app:signal, app:noise,
% app:priors, app:mcmc (labels of the appendix sections).

\subsection{Signal Model and Instrument Response}
\label{sec:problem-signal}
% quasi-monochromatic chirping binaries, the FastGB-style response, TDI A/E/T;
% frequency windows and why a window holds a whole response (app:signal, fig:band)

\subsection{Noise Model and Likelihood}
\label{sec:problem-likelihood}
% stationary Gaussian noise, the TDI-1.5 PSD, the Whittle likelihood, the SNR (app:noise)

\subsection{Priors and the Simulated Dataset}
\label{sec:problem-dataset}
% the per-window prior: log-uniform f0, Korol+22 chirp mass, SNR-defined amplitude,
% isotropic angles; four sources per window; the XS/S/B rungs (app:signal, app:priors; fig:priors)

\subsection{Bayesian Inference and the Global Fit}
\label{sec:problem-bayes}
% the posterior over an unordered set; exact symmetries and degeneracies; MCMC global
% fits (RJMCMC, parallel tempering) and their cost; what amortisation buys (app:mcmc)

% =======================================================================================
\section{Flow Matching on the Source Manifold}
\label{sec:method}
% =======================================================================================
% Drafted material: app:geometry, app:fm, app:sensitivity, app:conditioning,
% app:architecture, app:training.

\subsection{Conditional Flow Matching from Prior to Posterior}
\label{sec:method-cfm}
% the base distribution is the prior; geodesic paths; the objective (app:fm)

\subsection{A Product Manifold for a Set of Binaries}
\label{sec:method-geometry}
% chart, spheres and circles, the permutation-aligned coupling (app:geometry)

\subsection{Emphasising the End of the Path}
\label{sec:method-warp}
% the loss is blind to narrow posteriors; the warped clock s = 1 - (1 - t)^3 (app:sensitivity, fig:warp)

\subsection{Network Architecture}
\label{sec:method-architecture}
% conditioning input (app:conditioning, fig:window); MMDiT with a permutation-equivariant parameter stream
% (app:architecture, fig:architecture)

\subsection{Training and Sampling}
\label{sec:method-training}
% auxiliary heads, Muon, warmup-stable-decay cooldown, RK4 on the manifold (app:fm, app:training; alg:train, alg:sample)

% =======================================================================================
\section{Experiments}
\label{sec:experiments}
% =======================================================================================
% Drafted material: app:evaluation, app:xs, app:b, app:mcmc.

\subsection{Setup and Metrics}
\label{sec:experiments-setup}
% XS and B; 840 scored sources; width, found, offset, coverage, rank; Fisher reference
% (app:evaluation)

\subsection{Results on XS}
\label{sec:experiments-xs}
% the ablation ladder: warp, width, cooldown; detection and calibration (app:xs;
% fig:width-snr, fig:ladder, fig:calibration, fig:bias, fig:corner-*; tab:xs-ladder)

\subsection{The Whole Band}
\label{sec:experiments-b}
% B1: placeholders until the run finishes (app:b)

\subsection{Comparison with MCMC}
\label{sec:experiments-mcmc}
% validation of the forward model against Strub et al. (2022); flow against a
% parallel-tempered MCMC on selected injections: placeholders (app:mcmc)

\subsection{Cost}
\label{sec:experiments-cost}
% training cost per rung; sampling cost per window against MCMC (app:training, app:mcmc)

% =======================================================================================
\section{Conclusions}
\label{sec:conclusions}
% =======================================================================================
% Planned: what works (centred, calibrated-conservative posteriors that capture the
% degeneracies; the recipe); the remaining f0 floor (25x the ideal width at high SNR);
% limitations (fixed number of sources per window, no confusion foreground, windowing,
% bf16 resolution of the inputs on wide windows); next steps (MCMC validation, B,
% stitching windows into a global fit, a variable number of sources).


\section*{Software and Data}
% TODO: anonymous code link for review (the canna package: problem, network, training,
% scorecard, and the paper/scripts that make every figure and table).

% Acknowledgements should only appear in the accepted version.
% \section*{Acknowledgements}

\section*{Impact Statement}

This paper presents work whose goal is to advance the field of Machine Learning. There
are many potential societal consequences of our work, none which we feel must be
specifically highlighted here.

\bibliography{references}
\bibliographystyle{icml2026}

%%%%%%%%%%%%%%%%%%%%%%%%%%%%%%%%%%%%%%%%%%%%%%%%%%%%%%%%%%%%%%%%%%%%%%%%%%%%%%%%%%%%%%%%%%
% APPENDIX: all the drafted material, mathematical details first
%%%%%%%%%%%%%%%%%%%%%%%%%%%%%%%%%%%%%%%%%%%%%%%%%%%%%%%%%%%%%%%%%%%%%%%%%%%%%%%%%%%%%%%%%%
\newpage
\appendix
\onecolumn

% =======================================================================================
\section{Notation and Overview of the Appendices}
\label{app:notation}
% =======================================================================================

The appendices hold the complete technical description of the problem, the method and the
experiments. \Cref{tab:notation} collects the notation. The appendices are ordered from the
physics to the results:
\begin{itemize}
  \item \cref{app:signal}: the galactic-binary signal, the LISA response and the frequency
        windows that define a training example;
  \item \cref{app:noise}: the noise model, the likelihood, the signal-to-noise ratio, and how
        the forward model was validated;
  \item \cref{app:priors}: the prior over the sources in a window, its exact symmetries, and the
        simulated dataset;
  \item \cref{app:geometry}: the flow coordinates and the product manifold of an unordered set
        of sources;
  \item \cref{app:fm}: the conditional flow-matching objective, the warped clock, the
        auxiliary heads and the sampler;
  \item \cref{app:sensitivity}: why the flow-matching loss is nearly blind to the width of a
        narrow posterior, and what the warped clock changes;
  \item \cref{app:conditioning,app:architecture,app:training}: the conditioning input, the
        network, and the optimisation and compute;
  \item \cref{app:evaluation}: the evaluation protocol and the metrics;
  \item \cref{app:xs,app:b,app:mcmc}: the experiments on the XS rung, on the whole band (B),
        and the comparison with MCMC.
\end{itemize}

\begin{table}[h]
\caption{Notation.}
\label{tab:notation}
\centering
\small
\setlength{\tabcolsep}{4pt}
\begin{tabular}{lll}
\toprule
symbol & meaning & value or range \\
\midrule
$\Tobs$ & observation time & $2\,\text{yr} = 6.31\times10^{7}$\,s \\
bin & frequency resolution $1/\Tobs$ & $15.84$\,nHz \\
$S$ & number of sources in a window & 4 \\
$\theta = (f_0, \Mc, \Amp, \alpha, \delta, \psi, \iota, \phi_0)$ & one source; $(\alpha,\delta)$: RA, dec & \cref{tab:prior} \\
$\fdot(\Mc, f_0)$ & frequency derivative, derived from $(\Mc, f_0)$ & \cref{eq:fdot} \\
$R$ & response points: bins of one source's response & 16 (XS), 64 (S), 1024 (B) \\
$L$ & light travel time along an arm (equal arms) & $2.5\,\text{Gm}/c = 8.34$\,s \\
$K = 4R$ & bins in a window & 64, 256, 4096 \\
$k_0$, $\fw$ & first bin and centre frequency of a window & \cref{app:signal-windows} \\
$\tilde d_{k,c}$ & data in bin $k$ and TDI channel $c\in\{A,E,T\}$ & complex \\
$S_c(f)$ & one-sided noise PSD of channel $c$ (fractional frequency) & \cref{eq:psd} \\
$\inner{a}{b}$ & noise-weighted inner product & \cref{eq:inner} \\
$\snr$ & optimal matched-filter SNR, $\snr^2 = \inner{h}{h}$ & 7 to 1000 (prior) \\
$\pi_w$ & prior over the $S$ sources of window $w$ & \cref{app:priors} \\
$\varphi_w$ & chart from physical parameters to flow coordinates & \cref{eq:chart} \\
$\manifold$ & flow manifold of one source, embedded in $\R^{11}$ & $[-1,1]\times\R\times[-1,1]\times\Sph^2\times\Sph^2\times\Sph^1$ \\
$X = (x_1,\dots,x_S)$ & a set of sources in flow coordinates & $\manifold^S/\mathfrak{S}_S$ \\
$y$ & conditioning input (whitened, transformed data) & $2\times K\times 3$ \\
$t$, $s(t)$ & flow time and path position, $s = 1-(1-t)^p$ & $p=3$ (warped), $p=1$ (uniform) \\
$\vnet$ & the network's velocity field, weights $\eta$ & \cref{app:architecture} \\
$\Exp_x$, $\Log_x$ & exponential and logarithm maps & \cref{app:geometry} \\
$y_{ij}$ & response of the link received on spacecraft $i$ from $j$ & \cref{eq:link-response} \\
\bottomrule
\end{tabular}
\end{table}

% =======================================================================================
\section{The Galactic-Binary Signal and the LISA Response}
\label{app:signal}
% =======================================================================================

We compute the signal with JaxGB \citep{jaxgb}. This section reproduces its model in its own
notation and conventions, which are also those of the time-domain response simulator LISA GW
Response \citep{lisagwresponse}: the two codes share the triad, the polarisation convention
and the link response, and pass the same eight parameters to each other unchanged. Only
the heterodyned evaluation (\cref{app:signal-heterodyne}) and the TDI combination
(\cref{app:signal-tdi}) are specific to JaxGB.

\subsection{Waveform}
\label{app:signal-waveform}

A detached compact binary in the Galaxy, most often a pair of white dwarfs, emits a
quasi-monochromatic gravitational wave at twice its orbital frequency
\citep{nelemans2001gravitational,korol2017prospects}, which drifts slowly through radiation
reaction. The transverse-traceless strain of the plane wave is
\begin{equation}
  h_{ij}(\xi) = h_+(\xi)\,e^+_{ij} + h_\times(\xi)\,e^\times_{ij},\qquad
  \xi = t - \frac{\hat{\mathbf k}\cdot\mathbf x(t)}{c},
  \label{eq:xi}
\end{equation}
with $\hat{\mathbf k}$ the direction of propagation, $\mathbf x$ the position and
$e^{+,\times}_{ij}$ the polarisation tensors of the source frame. For a mildly chirping,
circular binary the two polarisations are
\begin{equation}
  h_+(\xi) = -\Amp\,\bigl(1+\cos^2\iota\bigr)\cos\Psi(\xi),\qquad
  h_\times(\xi) = -2\Amp\cos\iota\,\sin\Psi(\xi),\qquad
  \Psi(\xi) = -\phi_0 + 2\pi f_0\,\xi + \pi\fdot\,\xi^2,
  \label{eq:waveform}
\end{equation}
with strain amplitude $\Amp$, inclination $\iota$ of the orbital angular momentum to the line
of sight, initial phase $\phi_0$, frequency $f_0$ and frequency derivative $\fdot$, all
referred to the start of the observation. The sky position, right ascension $\alpha$ and
declination $\delta$ in the equatorial (ICRS) frame of the orbit files, defines the
orthonormal triad
\begin{equation}
  \hat{\mathbf u} = (\sin\alpha,\,-\cos\alpha,\,0),\qquad
  \hat{\mathbf v} = (-\sin\delta\cos\alpha,\,-\sin\delta\sin\alpha,\,\cos\delta),\qquad
  \hat{\mathbf k} = (-\cos\delta\cos\alpha,\,-\cos\delta\sin\alpha,\,-\sin\delta),
  \label{eq:triad}
\end{equation}
and with it the polarisation tensors of the solar-system barycentre (SSB),
$\epsilon^+_{ij} = (\hat{\mathbf u}\otimes\hat{\mathbf u} - \hat{\mathbf v}\otimes\hat{\mathbf v})_{ij}$
and
$\epsilon^\times_{ij} = (\hat{\mathbf u}\otimes\hat{\mathbf v} + \hat{\mathbf v}\otimes\hat{\mathbf u})_{ij}$.
The source frame is rotated from the SSB frame by the polarisation angle $\psi$,
$e^+_{ij} = \cos2\psi\,\epsilon^+_{ij} + \sin2\psi\,\epsilon^\times_{ij}$ and
$e^\times_{ij} = -\sin2\psi\,\epsilon^+_{ij} + \cos2\psi\,\epsilon^\times_{ij}$, so that
\begin{equation}
  h_{ij}(\xi) = h_+^{\rm SSB}(\xi)\,\epsilon^+_{ij} + h_\times^{\rm SSB}(\xi)\,\epsilon^\times_{ij},
  \qquad
  h_+^{\rm SSB} = h_+\cos2\psi - h_\times\sin2\psi,\qquad
  h_\times^{\rm SSB} = h_+\sin2\psi + h_\times\cos2\psi .
  \label{eq:pol-ssb}
\end{equation}
In complex form, $h_{ij}(\xi) = \mathrm{Re}\bigl\{\bigl(A_+^{\rm SSB}\epsilon^+_{ij}
+ A_\times^{\rm SSB}\epsilon^\times_{ij}\bigr)e^{\ii\Psi(\xi)}\bigr\}$ with
\begin{equation}
  A_+^{\rm SSB} = -\Amp\bigl[(1+\cos^2\iota)\cos2\psi + 2\ii\cos\iota\sin2\psi\bigr],\qquad
  A_\times^{\rm SSB} = -\Amp\bigl[(1+\cos^2\iota)\sin2\psi - 2\ii\cos\iota\cos2\psi\bigr].
  \label{eq:complex-amp}
\end{equation}
In LISA GW Response the same triad reads $\hat{\mathbf k} = -\mathbf e_r$,
$\hat{\mathbf u} = -\mathbf e_\phi$, $\hat{\mathbf v} = -\mathbf e_\theta$ in the spherical basis
of the ICRS, and its galactic-binary strain is \cref{eq:waveform,eq:pol-ssb} term by term.

For a circular binary driven by gravitational radiation alone, the frequency derivative is
fixed by the chirp mass $\Mc$,
\begin{equation}
  \fdot = \frac{96}{5}\,\pi^{8/3}\left(\frac{G\Mc}{c^3}\right)^{5/3} f_0^{11/3},
  \label{eq:fdot}
\end{equation}
and the amplitude by the distance $D$, $\Amp = 2(G\Mc)^{5/3}(\pi f_0)^{2/3}/(c^4 D)$. We
parametrise each source by $\theta = (f_0, \Mc, \Amp, \alpha, \delta, \psi, \iota, \phi_0)$,
derive $\fdot$ from \cref{eq:fdot}, and do not sample the distance: the amplitude prior is
set by the signal-to-noise ratio instead (\cref{app:noise-snr}). Tides and mass transfer,
which change $\fdot$ for some systems, are neglected.

\subsection{Single-Link Response}
\label{app:signal-response}

LISA measures the wave as relative frequency fluctuations of six laser links between three
spacecraft. Let $i$ be the receiving and $j$ the emitting spacecraft of link $ij$, $t_i$ the
time of reception and $t_j = t_i - L_{ij}(t_i)$ the time of emission, with $L_{ij}$ the light
travel time. The link response is \citep{finn2009response,cornish2009alternative}
\begin{equation}
  y_{ij}(t_i) = -\frac12\,\frac{H_{ij}(\xi_i) - H_{ij}(\xi_j)}{1 - \hat{\mathbf k}\cdot\hat{\mathbf x}_{ij}},
  \qquad
  H_{ij}(\xi) = \hat x^k_{ij}\hat x^l_{ij}\,h_{kl}(\xi)
  = F^+_{ij}\,h_+^{\rm SSB}(\xi) + F^\times_{ij}\,h_\times^{\rm SSB}(\xi),
  \label{eq:link-response}
\end{equation}
with $\xi_i = t_i - \hat{\mathbf k}\cdot\mathbf x_i(t_i)/c$,
$\xi_j = t_j - \hat{\mathbf k}\cdot\mathbf x_j(t_j)/c$, the unit vector
$\hat{\mathbf x}_{ij} = (\mathbf x_i(t_i) - \mathbf x_j(t_j))/\lVert\mathbf x_i(t_i) - \mathbf x_j(t_j)\rVert$
from the emitter to the receiver, and the antenna patterns
$F^+_{ij} = (\hat{\mathbf u}\cdot\hat{\mathbf x}_{ij})^2 - (\hat{\mathbf v}\cdot\hat{\mathbf x}_{ij})^2$
and $F^\times_{ij} = 2(\hat{\mathbf u}\cdot\hat{\mathbf x}_{ij})(\hat{\mathbf v}\cdot\hat{\mathbf x}_{ij})$.
This is the response of LISA GW Response, $y_{ij}(t) = [H_{ij}(t - L_{ij} - \hat{\mathbf k}\cdot\mathbf x_j/c)
- H_{ij}(t - \hat{\mathbf k}\cdot\mathbf x_i/c)]/[2(1 - \hat{\mathbf k}\cdot\mathbf n_{ij})]$,
with $\mathbf n_{ij}\equiv\hat{\mathbf x}_{ij}$.

For the galactic binary, define
$\mathcal A_{ij} = A_+^{\rm SSB}F^+_{ij} + A_\times^{\rm SSB}F^\times_{ij}$ and
\begin{equation}
  \alpha_{ij}(t_i) = \pi L_{ij}(t_i)\,\bigl(1 - \hat{\mathbf k}\cdot\hat{\mathbf x}_{ij}\bigr)\,
  \bigl(f_0 + \fdot\,\xi_i\bigr).
  \label{eq:alpha}
\end{equation}
With $L_{ij}\simeq\lVert\mathbf x_i(t_i) - \mathbf x_j(t_j)\rVert/c$ and the adiabatic
approximation $\fdot(\xi_i - \xi_j)/2\ll f_0$ (no significant change of frequency during a
light travel time), half the phase difference between emission and reception is
$\alpha_{ij}$, and the response becomes
\begin{equation}
  y_{ij}(t_i) = -\mathrm{Re}\left\{\frac{\ii\,\mathcal A_{ij}\sin\alpha_{ij}}
  {1 - \hat{\mathbf k}\cdot\hat{\mathbf x}_{ij}}\,
  e^{-\ii\left(\alpha_{ij} - \Psi(\xi_i)\right)}\right\}.
  \label{eq:link-gb}
\end{equation}

\subsection{Heterodyned Frequency-Domain Response}
\label{app:signal-heterodyne}

With $q = \operatorname{round}(f_0\Tobs)$ the frequency bin closest to $f_0$, the response
separates into a fast carrier and a slowly varying envelope,
$y_{ij}(t_i) = \mathrm{Re}\bigl\{\mathcal G_{ij}(t_i)\,e^{2\pi\ii qt_i/\Tobs}\bigr\}$, with
\begin{equation}
  \mathcal G_{ij}(t_i) = -\frac{\ii}{2}\,\frac{f_0 + \fdot\,\xi_i}{f_{\star}}\,\mathcal A_{ij}\,
  \operatorname{sinc}\alpha_{ij}\;
  \exp\Bigl\{-\ii\Bigl[\alpha_{ij} + 2\pi f_0\frac{\hat{\mathbf k}\cdot\mathbf x_i(t_i)}{c}
  - \Bigl(-\phi_0 + 2\pi\bigl(f_0 - \tfrac{q}{\Tobs}\bigr)t_i + \pi\fdot\,\xi_i^2\Bigr)\Bigr]\Bigr\},
  \label{eq:slow}
\end{equation}
where $\operatorname{sinc}x = \sin x/x$ and $f_\star = (2\pi L)^{-1}$, with $L$ in seconds.
The envelope varies on the scale of LISA's orbit, so its spectrum is confined to a few bins
around zero: the Doppler sideband $\pm f_0v_\oplus/c$ ($v_\oplus = 29.8$\,km/s, so
$\sim10^{-4}f_0$) plus the chirp drift $\fdot\,\Tobs$. JaxGB samples it at $R$ times
$t_m = m\Tobs/R$ and Fourier-transforms it, which gives the one-sided spectrum on the $R$ bins
starting at $k_{\min} = q - R/2$:
\begin{equation}
  \tilde y_{ij}\bigl(f_{q+k}\bigr) = \frac{\Tobs}{2R}\sum_{m=0}^{R-1}\mathcal G_{ij}(t_m)\,
  e^{-2\pi\ii km/R},\qquad k = -\tfrac R2,\dots,\tfrac R2 - 1,\qquad f_k = \frac{k}{\Tobs}.
  \label{eq:heterodyne}
\end{equation}
The factor $1/2$ is the positive-frequency half of the real signal, and the normalisation is
that of the continuous transform, $\tilde h(f) \simeq \int_0^{\Tobs}h(t)\,e^{-2\pi\ii ft}\,\dd t$.
This is the fast galactic-binary response of \citet{cornish2007tests}, vectorised over
sources and differentiable. The baseline JaxGB model makes two further approximations,
accurate at the percent level: the emitter's position is taken at the time of reception,
$\mathbf x_j(t_j)\simeq\mathbf x_j(t_i)$, so that $\hat{\mathbf x}_{ji} = -\hat{\mathbf x}_{ij}$
and $\mathcal A_{ji} = \mathcal A_{ij}$, and all six light travel times are equal and constant,
$L_{ij} = L$, their average. With the equal-arm orbits of LISA Orbits \citep{lisaorbits},
$L = 2.5\,\text{Gm}/c = 8.34$\,s and $f_\star = 19.1$\,mHz.

\subsection{Time-Delay Interferometry}
\label{app:signal-tdi}

Laser frequency noise is cancelled by time-delay interferometry
\citep[TDI;][]{tinto2021time}. JaxGB forms the first-generation (1.5) Michelson
combinations for equal and constant arms on the envelopes of \cref{eq:slow}, with the
instantaneous frequency $f = f_0 + \fdot\,t$ and the delay phase $x = 2\pi fL$:
\begin{equation}
  X = -2\ii\sin x\;e^{-\ii x}\Bigl[\bigl(\mathcal G_{21} - \mathcal G_{31}\bigr)
  + e^{-\ii x}\bigl(\mathcal G_{12} - \mathcal G_{13}\bigr)\Bigr],
  \label{eq:tdi-x}
\end{equation}
with $Y$ and $Z$ from the cyclic permutation $1\to2\to3\to1$, Fourier-transformed as in
\cref{eq:heterodyne}. Since $-2\ii\sin x\,e^{-\ii x} = -(1 - D^2)$ with $D = e^{-\ii x}$ the
delay operator, \cref{eq:tdi-x} is $(1-D^2)$ times a combination of the four links that
involve spacecraft 1. The standard equal-arm TDI-1.5 Michelson combination
\citep[eq.~14]{babak2021lisa} pairs the delays the other way round,
$(1-D^2)\bigl[(y_{13}-y_{12}) + D\,(y_{31}-y_{21})\bigr]$; the two differ by
$(1-D^2)(1-D)\bigl[(y_{13}-y_{31}) - (y_{12}-y_{21})\bigr]$, which is of higher order in $x$.
On JaxGB's own link responses, for sources drawn from our prior, they agree to
$2.4\times10^{-4}$ (relative) up to 3\,mHz and to $7\times10^{-3}$ at 10\,mHz, with powers
equal to 0.2\%. We use the quasi-orthogonal channels \citep{prince2002lisa}
\begin{equation}
  A = \frac{Z-X}{\sqrt2},\qquad E = \frac{X-2Y+Z}{\sqrt6},\qquad T = \frac{X+Y+Z}{\sqrt3},
  \label{eq:aet}
\end{equation}
whose noises are uncorrelated for equal arms (\cref{app:noise-psd}). The isotropic sky prior
does not depend on the frame; sources quoted in ecliptic coordinates, as in the LISA Data
Challenges, must be rotated to the equatorial frame of the orbits.

\subsection{Frequency Windows}
\label{app:signal-windows}

A training example is one window of $K=4R$ consecutive frequency bins, starting at an even
bin $k_0$, with centre frequency $\fw = (k_0 + K/2)/\Tobs$. The window holds $S=4$ sources,
whose $R$-bin responses are added at their own offsets $k_{\min}-k_0$. A guard band of $R/2$
bins at each edge keeps every response inside the window: sources are drawn with
\begin{equation}
  f_0 \in \bigl[f_{\rm lo}(w), f_{\rm hi}(w)\bigr]
  = \Bigl[\max\bigl(\tfrac{k_0 + R/2}{\Tobs}, f_{\min}\bigr),\;
          \min\bigl(\tfrac{k_0 + K - R/2}{\Tobs}, f_{\max}\bigr)\Bigr],
  \label{eq:window-interior}
\end{equation}
an interior of $3R$ bins. The window is placed by drawing $f_c$ log-uniformly on
$[f_{\min}, f_{\max}]$ and snapping the start to the nearest even bin
$k_0 = 2\operatorname{round}(f_c\Tobs/2 - R)$, clipped to the valid range. Windows therefore
tile the band with density $\propto 1/f$, and $f_0$ is log-uniform over the band to within
the snapping.

\paragraph{How many bins a source needs.}
Half of a source's response window must hold the Doppler sideband plus the whole chirp
drift, so a source needs
\begin{equation}
  n_{\rm span}(f_0,\Mc) = 2\left\lceil\Bigl(\frac{f_0 v_\oplus}{c}
  + \fdot(\Mc, f_0)\,\Tobs\Bigr)\Tobs\right\rceil \le R
  \label{eq:span}
\end{equation}
bins. For a given $R$, the highest frequency the prior can reach is set by the heaviest
binary, $\Mc_{\max} = 1.4\cdot2^{-1/5} = 1.219\,\Msun$, and solves
$(v_\oplus/c)\,f_0 + \kappa\,\Mc_{\max}^{5/3}\Tobs f_0^{11/3} = R/(2\Tobs)$, with
$\kappa$ the constant of \cref{eq:fdot}. In $u = f_0^{1/3}$ this
is a degree-11 trinomial, solved numerically by Newton's method from the smaller of the
Doppler-only root $Rc/(2\Tobs v_\oplus)$ and the chirp-only root. The two terms scale very
differently: the Doppler term is linear in $f_0$ and the chirp term goes as $f_0^{11/3}$, so
once the chirp dominates (above $\sim4$\,mHz for the heaviest binaries,
\cref{fig:band}c) quadrupling $R$ buys only $4^{3/11}\simeq1.46$ times more bandwidth.

\paragraph{Rungs.}
The problem comes in three sizes that differ only in $R$ and hence in how much of the band a
window covers (\cref{tab:rungs}). With the full band, $R$ is sized automatically as the
power of two above $n_{\rm span}(f_{\max},\Mc_{\max})$, with a floor of 256 points: at
12\,mHz the Doppler sideband is 75 bins and the chirp of the heaviest binary 291 bins, so
$n_{\rm span}=734$ and $R=1024$, which could hold the heaviest binaries up to 13.3\,mHz.

\begin{table}[h]
\caption{The three problem sizes. They share the priors, the four sources per window, the
network and the training recipe; only the response points $R$, and with them the window and
the number of data tokens, differ. The band fraction is of $\log f$ over 0.1--12\,mHz.}
\label{tab:rungs}
\centering
\small
\begin{tabular}{lrrrrll}
\toprule
rung & $R$ & window $K$ [bins] & interior [bins] & data tokens & $f_0$ range & band \\
\midrule
XS & 16 & 64 & 48 & 32 & 0.1--1.26\,mHz & 53\% \\
S & 64 & 256 & 192 & 128 & 0.1--4.15\,mHz & 78\% \\
B & 1024 & 4096 & 3072 & 2048 & 0.1--12\,mHz & 100\% \\
\bottomrule
\end{tabular}
\end{table}

\begin{figure}[h]
  \centering
  \includegraphics[width=\textwidth]{figures/band.pdf}
  \caption{The band as the training prior sees it. (a) The TDI-1.5 noise PSDs of the $A$,
  $E$ and $T$ channels in fractional-frequency units, which whiten the data
  (\cref{eq:psd}). (b) The amplitude prior: the band between a sky-averaged SNR of 7 and of
  1000 at each frequency (\cref{eq:amp-window}); the dot is the faint source of
  \citet{strub2022bayesian} at SNR 7.93, used to validate the likelihood (\cref{app:mcmc}).
  (c) The half-span of a source's response, Doppler sideband plus chirp drift over
  $\Tobs$ (\cref{eq:span}), against the half-window of each rung; the dots mark the highest
  frequency each rung can hold for the heaviest binary in the prior: 1.26, 4.15 and
  13.3\,mHz (the prior itself stops at 12\,mHz).}
  \label{fig:band}
\end{figure}

\subsection{What the Simulator Leaves Out}
\label{app:signal-limits}

The training distribution is deliberately simpler than real LISA data. Every window holds
exactly four resolvable binaries and nothing else: there is no confusion foreground from the
millions of unresolved binaries \citep{timpano2006characterizing,karnesis2021characterization},
no other source class, no gaps or glitches, and the noise is stationary and known. Sources in
neighbouring windows never leak in, because the guard band keeps every response inside its
own window. For $R<256$ (XS and S) the slow response is sampled critically rather than
generously, which leaves per-bin amplitudes of the highest-frequency XS sources with a
$\sim$10--25\% aliasing error against a 1024-point reference while preserving the total
power; the same simulator generates the training data and the test injections, so this
affects realism, not consistency. Using the posteriors inside a global fit would require
lifting the first two simplifications: a variable number of sources per window, and sources
that straddle window edges.

% =======================================================================================
\section{Noise Model, Likelihood and Signal-to-Noise Ratio}
\label{app:noise}
% =======================================================================================

\subsection{Instrument Noise}
\label{app:noise-psd}

We use the two noise terms of the LISA science requirements as given by
\citet{babak2021lisa}: the optical metrology system (OMS) and the test-mass acceleration
noise of one link, as displacement and acceleration spectra (their eqs.~9 and 11),
\begin{align}
  P_{\rm oms}(f) &= (15\,\text{pm})^2\left[1+\Bigl(\frac{2\,\text{mHz}}{f}\Bigr)^4\right]
  \,\text{m}^2\,\text{Hz}^{-1}, \nonumber\\
  P_{\rm acc}(f) &= \Bigl(3\,\frac{\text{fm}}{\text{s}^2}\Bigr)^2
  \left[1+\Bigl(\frac{0.4\,\text{mHz}}{f}\Bigr)^2\right]
  \left[1+\Bigl(\frac{f}{8\,\text{mHz}}\Bigr)^4\right]
  \,\text{m}^2\,\text{s}^{-4}\,\text{Hz}^{-1}.
  \label{eq:single-link}
\end{align}
A link measures a Doppler shift, so in the relative-frequency units of the response
(\cref{eq:link-response}) a displacement noise enters as $(2\pi f/c)^2$ in power and an
acceleration noise picks up a further $(2\pi f)^{-4}$ (their eqs.~10 and 13):
\begin{equation}
  S_{\rm oms}(f) = \Bigl(\frac{2\pi f}{c}\Bigr)^2P_{\rm oms}(f),\qquad
  S_{\rm acc}(f) = \frac{P_{\rm acc}(f)}{(2\pi f)^2\,c^2}.
\end{equation}
For equal arms, identical noise levels on all links and uncorrelated noises, the TDI-1.5
Michelson combination has the PSD $S_X$ of \citet[eq.~19]{babak2021lisa}, and $X$ and $Y$ have
the cross spectral density $S_{XY}$, which we obtain from their TDI-2.0 expression (eq.~129)
by removing the generation-2 transfer factor $4\sin^2 2x$:
\begin{equation}
  S_X = 16\sin^2x\,\bigl[S_{\rm oms} + (3+\cos2x)\,S_{\rm acc}\bigr],\qquad
  S_{XY} = -8\sin^2x\,\cos x\,\bigl[S_{\rm oms} + 4S_{\rm acc}\bigr],\qquad
  x = 2\pi fL .
\end{equation}
The channels of \cref{eq:aet} diagonalise the $XYZ$ covariance, with $S_A = S_E = S_X - S_{XY}$
and $S_T = S_X + 2S_{XY}$ (their eqs.~127--128):
\begin{align}
  S_A(f) = S_E(f) &= 8\sin^2x\,\bigl[(2+\cos x)\,S_{\rm oms}
  + 2\,(3 + 2\cos x + \cos2x)\,S_{\rm acc}\bigr], \nonumber\\
  S_T(f) &= 16\sin^2x\,(1-\cos x)\,\bigl[S_{\rm oms} + 2\,(1-\cos x)\,S_{\rm acc}\bigr].
  \label{eq:psd}
\end{align}
$S_A$ is their eq.~130 divided by $4\sin^22x$. These are the PSDs our simulator draws the
noise from and our likelihood whitens with; we checked them against the expressions above to
machine precision, with the same $L = 8.34$\,s as the response. \Cref{fig:band}a shows them.
The units matter more than they seem: writing the noise as a fractional length (dividing by
the arm length rather than by $c/2\pi f$) leaves the whitening wrong by $x^2$, a tilt across
the band rather than a constant, which would silently re-weight low-frequency against
high-frequency sources. We checked the units by forming the sky-averaged sensitivity from
the Monte-Carlo-averaged response of the $A$ and $E$ channels and comparing it with the
analytic curve of \citet{robson2019construction}: the ratio is flat (log-log slope $-0.001$)
at the expected $1/2$ below 4\,mHz, against a slope of $-2$ for the fractional-length form.

\subsection{Data Model and Likelihood}
\label{app:noise-likelihood}

In a window, the data are the sum of the four responses and Gaussian noise,
\begin{equation}
  \tilde d_{k,c} = \sum_{s=1}^{S}\tilde h_c(f_k;\theta_s) + \tilde n_{k,c},\qquad
  \tilde n_{k,c} \sim \CN\!\left(0,\ \tfrac{1}{2}S_c(f_k)\,\Tobs\right),
  \label{eq:data}
\end{equation}
independent across bins $k$ and channels $c\in\{A,E,T\}$: the noise is stationary, and the
three channels are uncorrelated for equal arms. With the noise-weighted inner product
\begin{equation}
  \inner{a}{b} = 4\,\mathrm{Re}\sum_{c}\sum_{k}\frac{a_{k,c}\,b^*_{k,c}}{S_c(f_k)\,\Tobs},
  \label{eq:inner}
\end{equation}
the discrete form of $4\,\mathrm{Re}\int a\,b^*/S\,\dd f$, the log-likelihood of the window
is the Whittle likelihood
\begin{equation}
  \log p(d\mid\theta_{1:S}) = -\tfrac12\inner{d-h}{d-h} + \text{const}
  = -\sum_{k,c}\frac{\bigl|\tilde d_{k,c}-\tilde h_{k,c}(\theta_{1:S})\bigr|^2}{S_c(f_k)\,\Tobs/2}
  + \text{const}.
  \label{eq:loglik}
\end{equation}
Training never evaluates it, because the flow learns from simulated pairs alone, but it
defines the posterior that the flow must reproduce, and the Fisher reference
(\cref{app:evaluation-fisher}) and the MCMC baseline (\cref{app:mcmc}) use it directly.

\subsection{Signal-to-Noise Ratio and the Amplitude Prior}
\label{app:noise-snr}

The optimal matched-filter SNR is $\snr^2 = \inner{h}{h}$: for noiseless data $d = h$, it is
twice the log-likelihood gain of the true sources over no source,
$\snr^2 = 2\,[\log p(d\mid\theta_{1:S}) - \log p(d\mid\text{no source})]$. We quote two
versions: the SNR of one source (its response alone) and the SNR of a window (all four
sources). A fixed box in amplitude would be a very different prior at the two ends of the
band: the same $\Amp$ is undetectable at 0.1\,mHz and very loud at 5\,mHz. As in
\citet{strub2022bayesian}, we instead set the amplitude prior through the SNR. For a
monochromatic source averaged over sky position, polarisation, inclination and phase,
\begin{equation}
  \snr^2 \simeq C\,\frac{\Amp^2\,\Tobs}{S_n(f)},\qquad
  S_n(f) = \frac{10}{3\,(cL)^2}\left[P_{\rm oms} + \frac{4P_{\rm acc}}{(2\pi f)^4}\right]
  \left[1+0.6\Bigl(\frac{f}{f_\star}\Bigr)^2\right],\quad f_\star = \frac{1}{2\pi L},
  \label{eq:snr-sky}
\end{equation}
with $cL = 2.5$\,Gm the arm length. $S_n$ is the sky-averaged sensitivity of
\citet[eq.~13]{robson2019construction} in its low-frequency form: their acceleration term
carries $2\bigl(1+\cos^2(f/f_\star)\bigr)$, which tends to the 4 above for $f\ll f_\star$; the
two differ by at most 0.7\% in $S_n$ (0.3\% in amplitude) between 0.1 and 12\,mHz, because
the acceleration noise is negligible where they differ. Rather than take the constant $C$
from a paper whose normalisation conventions differ from ours, we measure it by Monte Carlo
against our own $\snr$ with this $S_n$: $C = 3.225$, flat to 0.1\% below 1.5\,mHz and
drifting by 7\% at 12\,mHz, where the $1+0.6(f/f_\star)^2$ approximation runs out. A window at
centre frequency $\fw$ then has the amplitude range
\begin{equation}
  \Amp \in \bigl[\Amp_{\rm lo}(\fw), \Amp_{\rm hi}(\fw)\bigr]
  = [\snr_{\min}, \snr_{\max}]\,\sqrt{\frac{S_n(\fw)}{C\,\Tobs}},
  \qquad [\snr_{\min},\snr_{\max}] = [7, 1000],
  \label{eq:amp-window}
\end{equation}
evaluated at the window centre rather than at each source's frequency: the sensitivity is
flat across a window, and this keeps the amplitude coordinate (\cref{eq:chart}) a function
of the window alone. \Cref{fig:band}b shows the range across the band and \cref{fig:priors}b
the per-source SNR it produces.

\subsection{Validation of the Forward Model}
\label{app:noise-validation}

Before training, we checked the simulator and the likelihood against an independent analysis
and against their own statistics.
\begin{itemize}
  \item \textbf{SNR against a published source.} For the faint source of
        \citet[Table VII]{strub2022bayesian} ($f_0 = 1.40457$\,mHz,
        $\log_{10}\fdot = -17.3$, $\Amp = 4.55\times10^{-23}$, ecliptic longitude and latitude
        $(1.5, 0.2)$ rotated to the equatorial $(\alpha,\delta)$ of the orbits,
        $\iota = 1.2$, $\psi = 2$, $\phi_0 = 3$, two years), our SNR is 7.951 against their
        7.93. Their value is the SNR in their noise realisation, a draw from
        $\Normal(\snr_{\rm opt}, 1)$, so the agreement (0.26\%) is better than the check can
        resolve.
  \item \textbf{Posterior against a published source.} A parallel-tempered ensemble MCMC on
        our likelihood (\cref{app:mcmc}) reproduces the width and shape of their posterior:
        the $f_0$--$\fdot$ correlation, the prior-limited $\fdot$, the $\cos\iota$--$\Amp$
        anticorrelation, and $\sigma_{f_0} = 1.40$\,nHz against their $+1.9/-1.5$\,nHz.
        Their narrower $\cos\iota$ comes from sampling inside a $\pm3\sigma$ Fisher box
        around the maximum-likelihood point, which truncates the tail towards face-on; inside
        the same box we reproduce their width. Over 48 noise realisations, the gain from the
        optimal SNR to the maximum-likelihood SNR matches theirs (0.32), which tests the
        curvature of the likelihood.
  \item \textbf{Noise generator.} Whitened by the same PSD the likelihood uses, 3000
        pure-noise windows give complex values of unit variance ($E|z|^2 = 1.0000$, variance
        0.5000 in each of the real and imaginary parts, correlation $-2\times10^{-4}$), an
        empirical PSD equal to $S_c$ to Monte-Carlo precision with no frequency tilt, and
        $\inner{n}{n}$ distributed as the expected $\chi^2$ (Kolmogorov--Smirnov $p = 0.83$).
  \item \textbf{Response and TDI conventions.} JaxGB's waveform and link response
        (\cref{eq:waveform,eq:pol-ssb,eq:link-response}) are those of the documentation of
        JaxGB and of LISA GW Response, which share the triad, the polarisation rotation and
        the link response term by term. JaxGB's TDI combination (\cref{eq:tdi-x}) agrees with
        the standard equal-arm TDI-1.5 $X$ of \citet[eq.~14]{babak2021lisa}, built from the
        same link responses, to $2.4\times10^{-4}$ up to 3\,mHz and $7\times10^{-3}$ at 10\,mHz
        (\cref{app:signal-tdi}).
  \item \textbf{Noise PSDs.} The $A$, $E$, $T$ PSDs of the simulator equal
        \cref{eq:psd}, that is eqs.~19 and 127--130 of \citet{babak2021lisa} at TDI 1.5, to
        machine precision.
  \item \textbf{Self-consistency.} The two forms of $\snr^2$ above agree to the last digit, and an MCMC on a noiseless injection recovers the truth at the
        centre of the posterior with $\hat R\approx1.000$ on all eight parameters.
\end{itemize}

% =======================================================================================
\section{Priors, Symmetries and the Simulated Dataset}
\label{app:priors}
% =======================================================================================

\subsection{The Prior over a Window}
\label{app:priors-table}

Given a window $w$ (\cref{app:signal-windows}), the $S=4$ sources are independent and
identically distributed under the per-source prior of \cref{tab:prior}, so the prior over
the window is $\pi_w(\theta_{1:S}) = \prod_{s}\pi_w(\theta_s)$. Only the frequency and the
amplitude depend on the window.

\begin{table}[h]
\caption{Prior of one source in window $w$. $\Unif$ is uniform; $f_{\rm lo,hi}$ and
$\Amp_{\rm lo,hi}$ are given by \cref{eq:window-interior,eq:amp-window}.}
\label{tab:prior}
\centering
\small
\begin{tabular}{llll}
\toprule
parameter & prior & support & note \\
\midrule
$f_0$ & $\log f_0\sim\Unif(\log f_{\rm lo}, \log f_{\rm hi})$ & window interior, $3R$ bins & window placed log-uniformly in the band \\
$\Mc$ & DWD population, \cref{eq:mc-prior} & $[0.131, 1.219]\,\Msun$ & \citet{korol2022observationally} \\
$\fdot$ & derived, \cref{eq:fdot} & & not sampled \\
$\Amp$ & $\log\Amp\sim\Unif(\log\Amp_{\rm lo}, \log\Amp_{\rm hi})$ & sky-averaged SNR 7 to 1000 & set at the window centre \\
$\alpha$ & $\Unif(0, 2\pi)$ & & isotropic sky (right ascension) \\
$\delta$ & $\sin\delta\sim\Unif(-1, 1)$ & & isotropic sky (declination) \\
$\psi$ & $\Unif(0, 2\pi)$ & & double cover, \cref{app:priors-symmetries} \\
$\iota$ & $\cos\iota\sim\Unif(-1, 1)$ & & isotropic orientation \\
$\phi_0$ & $\Unif(0, 2\pi)$ & & \\
\bottomrule
\end{tabular}
\end{table}

\paragraph{Frequency.}
The frequency is log-uniform inside the window, and windows are placed log-uniformly over the
band, so $f_0$ is log-uniform over 0.1\,mHz to $f_{\max}(R)$. This is a training choice, not a
population model: the Galactic population falls steeply with frequency
\citep{maoz2018separation,korol2022observationally}, and following it would starve the
high-frequency end of the band, where sources are rare but loud and well measured.

\paragraph{Chirp mass.}
We use the observationally driven double-white-dwarf population of
\citet{korol2022observationally}. The heavier component follows the single-white-dwarf mass
function of \citet{kepler2015new}, a three-component Gaussian mixture truncated to
$[m_{\min}, m_{\max}] = [0.15, 1.4]\,\Msun$, and the lighter one is uniform below it, the flat
mass-ratio distribution of close binaries \citep{moe2017mind}:
\begin{equation}
  m_1 \sim \sum_{j=1}^{3} w_j\,\Normal_{[m_{\min}, m_{\max}]}(\mu_j, \sigma_j^2),\qquad
  m_2\mid m_1 \sim \Unif(m_{\min}, m_1),\qquad
  \Mc = \frac{(m_1m_2)^{3/5}}{(m_1+m_2)^{1/5}},
  \label{eq:mc-prior}
\end{equation}
with $w = (0.81, 0.14, 0.05)$, $\mu = (0.65, 0.57, 0.81)\,\Msun$ and
$\sigma = (0.044, 0.097, 0.187)\,\Msun$. Truncation re-weights the components by the mass
they keep, $w_j\,[\Phi(b_j) - \Phi(a_j)]$, with $a_j, b_j$ the standardised bounds. The
induced chirp mass (\cref{fig:priors}a) is unimodal with a mode near $0.47\,\Msun$ and median
$0.43\,\Msun$; its support runs between the equal-mass extremes, $m\,2^{-1/5}$. The branch of
\citet{korol2022observationally} for extremely low-mass primaries ($m_1<0.25\,\Msun$) is never
reached by this mixture and is left out. This distribution peaks at higher chirp mass than
the population-synthesis catalogues used in the LISA Data Challenges \citep{korol2017prospects,baghi2022lisa}, which peak near $0.25\,\Msun$ \citep[Fig.~3]{korol2022observationally}.

\paragraph{Amplitude.}
The amplitude is log-uniform between a sky-averaged SNR of 7 and of 1000
(\cref{eq:amp-window}), the search range of \citet{strub2022bayesian}. Because sky position,
inclination and polarisation scatter the actual SNR around its sky average, the per-source
SNR (\cref{fig:priors}b) is log-uniform with soft edges: on XS, 3.9\% of the sources fall
below 7 and 1.8\% above 1000, with median 77 (log-uniform between 7 and 1000 gives
$\sqrt{7000} = 84$). The amplitude is drawn independently of the chirp mass and the distance;
drawing $\Amp$ from $(\Mc, f_0, D)$ with a Galactic distance model would make the prior
astrophysically consistent at the price of a much less controlled SNR distribution.

\begin{figure}[h]
  \centering
  \includegraphics[width=0.5\textwidth]{figures/priors.pdf}
  \caption{The two non-trivial priors. (a) Chirp mass of the observationally driven
  double-white-dwarf population (\cref{eq:mc-prior}), 200\,000 draws. (b) Optimal SNR of single
  sources under the XS prior, 8192 sources, against the log-uniform target between 7 and 1000.}
  \label{fig:priors}
\end{figure}

\subsection{Exact Symmetries}
\label{app:priors-symmetries}

The posterior inherits two groups of exact symmetries, and a sampler has to reproduce both.

\paragraph{Polarisation and phase.}
The source-frame polarisation tensors of \cref{eq:xi} are rotated from the SSB ones by
$2\psi$, so they have period $\pi$ in $\psi$, and a quarter turn flips their sign,
$e^{+,\times}(\psi+\pi/2) = -e^{+,\times}(\psi)$, as it flips $h_{+,\times}^{\rm SSB}$ in
\cref{eq:pol-ssb}. A half-cycle of the phase $\Psi$ undoes it, so the waveform
\cref{eq:waveform} is invariant under
\begin{equation}
  (\psi, \phi_0) \;\mapsto\; \bigl(\psi + k\tfrac{\pi}{2},\ \phi_0 + k\pi\bigr)
  \pmod{2\pi},\qquad k = 0,1,2,3.
  \label{eq:psi-phi-symmetry}
\end{equation}
Our prior covers $\psi\in[0,2\pi)$, twice the range that would remove the $k=2$ copy, so every
source's posterior has exactly four copies in $(\psi,\phi_0)$. We keep the double cover
because the orientation chart of \cref{app:geometry} needs $\psi$ to range over a full
circle; the flow learns all four modes (\cref{fig:corner-median}).

\paragraph{Exchangeable sources.}
The sources in a window are i.i.d.\ under the prior and enter the likelihood through their
sum, so prior, likelihood and posterior are invariant under the $S! = 24$ permutations of the
source labels: the label-switching problem of mixture models
\citep{stephens2000dealing,jasra2005markov}. Together with \cref{eq:psi-phi-symmetry}, every
mode of a window's posterior appears in $24\times4^4 = 6144$ exactly equivalent copies. We
handle the permutations in the geometry and the architecture (\cref{app:geometry-set,app:architecture-equivariance}), so that the flow works on unordered sets, and let the
network learn the four $(\psi,\phi_0)$ copies.

\paragraph{Approximate degeneracies.}
The data constrain some combinations much better than others: amplitude against inclination
(a face-on source at larger distance looks like an inclined one nearer by), the
$(\psi,\phi_0)$ combination that survives for nearly face-on sources, where the wave is close
to circularly polarised, and, at low frequency, where the chirp is not resolved
($\fdot\Tobs^2\ll1$ bin, \cref{fig:band}c), the chirp mass, whose posterior is the prior.

\subsection{The Simulated Dataset}
\label{app:priors-dataset}

There is no fixed dataset. Each training step draws, inside the compiled step, 256 fresh
windows, their sources from the prior, the responses, the noise, and the conditioning input,
so that a run of $10^6$ steps sees $2.6\times10^8$ independent windows and $10^9$ sources. On XS
a whole step costs 38.9\,ms on an A100; fitting the step time of XS and S against the
network's FLOPs attributes about 23\,ms to a fixed per-step cost, which includes the
simulator, and 15\,ms to the network (\cref{app:training-compute}). On a laptop GPU the
simulator measured below 1\% of the step. Training on fresh simulations removes overfitting as a
concern and turns every comparison between runs into a comparison of optimisation, not of
generalisation.

% =======================================================================================
\section{Flow Coordinates and the Source Manifold}
\label{app:geometry}
% =======================================================================================

\subsection{The Chart}
\label{app:geometry-chart}

The flow does not work in physical units. A chart $\varphi_w$, which depends on the window,
maps each source to eleven numbers on a product manifold:
\begin{align}
  x_f &= \frac{2\ln f_0 - \ln f_{\rm lo} - \ln f_{\rm hi}}{\ln f_{\rm hi} - \ln f_{\rm lo}}\in[-1,1], &
  x_{\Mc} &= \frac{\ln\Mc - \mu_{\Mc}}{\sigma_{\Mc}}\in\R, &
  x_{\Amp} &= \frac{2\ln\Amp - \ln\Amp_{\rm lo} - \ln\Amp_{\rm hi}}{\ln\Amp_{\rm hi} - \ln\Amp_{\rm lo}}\in[-1,1], \nonumber\\
  n_{\rm sky} &= \begin{pmatrix}\cos\delta\cos\alpha\\ \cos\delta\sin\alpha\\ \sin\delta\end{pmatrix}\in\Sph^2, &
  n_{\rm or} &= \begin{pmatrix}\sin\iota\cos\psi\\ \sin\iota\sin\psi\\ \cos\iota\end{pmatrix}\in\Sph^2, &
  e_{\phi} &= \begin{pmatrix}\cos\phi_0\\ \sin\phi_0\end{pmatrix}\in\Sph^1,
  \label{eq:chart}
\end{align}
so that a source lives on
$\manifold = [-1,1]\times\R\times[-1,1]\times\Sph^2\times\Sph^2\times\Sph^1$, of dimension 8,
embedded in $\R^{11}$. Frequency and amplitude are whitened relative to the window
(\cref{eq:window-interior,eq:amp-window}), so the network always predicts quantities of order
one, and the absolute frequency reaches it through the conditioning instead
(\cref{app:architecture}). The log chirp mass is standardised by its prior mean and standard
deviation, computed once by quadrature over $(m_1, m_2/m_1)$ on a $256\times256$ grid. The
inverse chart uses $\operatorname{atan2}$ for every angle.

The chart is chosen so that the prior is simple on $\manifold$: $x_f$ and $x_{\Amp}$ are
uniform on $[-1,1]$, $n_{\rm sky}$ and $n_{\rm or}$ are uniform on their spheres (isotropic
sky, $\psi$ uniform and $\cos\iota$ uniform), and $e_\phi$ is uniform on the circle. Only
$x_{\Mc}$ keeps a non-trivial density. Angles then have no branch cuts and the sky no poles:
a posterior that straddles $\alpha = 0$ or $\phi_0 = 0$ is a single connected mode on the
manifold.

\subsection{Geodesics}
\label{app:geometry-maps}

Each factor comes with its exponential map $\Exp_x(v)$, which follows the geodesic from $x$
with initial velocity $v$ for unit time, and its logarithm $\Log_x(x')$, the initial velocity
of the geodesic from $x$ to $x'$:
\begin{itemize}
  \item \textbf{line} ($x_{\Mc}$): $\Log_x(x') = x'-x$, $\Exp_x(v) = x+v$;
  \item \textbf{interval} ($x_f$, $x_{\Amp}$): the same, with $\Exp$ clipped to $[-1,1]$. The
        segment between two interior points stays inside, so the clip only guards the
        sampler against overshooting;
  \item \textbf{sphere} ($\Sph^{n-1}$, unit radius): with
        $\vartheta = \operatorname{atan2}\bigl(\lVert x' - \langle x,x'\rangle x\rVert,
        \langle x,x'\rangle\bigr)$ the angle between the points,
        \begin{equation}
          \Log_x(x') = \vartheta\,\frac{x' - \cos\vartheta\,x}{\lVert x' - \cos\vartheta\,x\rVert},
          \qquad
          \Exp_x(v) = \cos\lVert v_\perp\rVert\,x + \frac{\sin\lVert v_\perp\rVert}{\lVert v_\perp\rVert}\,v_\perp,
          \quad v_\perp = v - \langle v, x\rangle x,
        \end{equation}
        with an arbitrary direction orthogonal to $x$ at the antipode, where the geodesic is
        not unique. $\Exp$ discards the radial part of $v$, which keeps the sampler on the
        sphere;
  \item \textbf{product} $\manifold$: factor by factor, with the product metric
        $d(x,x')^2 = \sum_{\rm factors}\lVert\Log_x(x')\rVert^2$.
\end{itemize}
The geodesic interpolation between two points is
$\gamma(s) = \Exp_{x_0}\bigl(s\Log_{x_0}(x_1)\bigr)$, $s\in[0,1]$, with constant speed; we
obtain its velocity $\dot\gamma(s)$ by automatic differentiation of $\gamma$ with respect to $s$.

\subsection{An Unordered Set of Sources}
\label{app:geometry-set}

A window holds a set of $S$ sources, a point of the quotient $\manifold^S/\mathfrak S_S$ by the
permutation group. We work with an ordered representative and choose the order when it
matters, which is when a path is drawn between two sets. For base and target sets $X_0$ and
$X_1$, we align the target to the base,
\begin{equation}
  \sigma^\star = \operatorname*{arg\,min}_{\sigma\in\mathfrak S_S}
  \sum_{i=1}^{S} d\bigl(x_{0,i},\, x_{1,\sigma(i)}\bigr)^2,
  \qquad
  \Log_{X_0}(X_1) = \Bigl(\Log_{x_{0,i}}\bigl(x_{1,\sigma^\star(i)}\bigr)\Bigr)_{i=1}^{S},
  \label{eq:set-log}
\end{equation}
by exhaustive search over the $S!=24$ permutations (for larger sets the implementation falls
back to randomised pairwise swaps). The aligned path never sends two slots across each other
when swapping their targets would shorten it. \Cref{eq:set-log} is the logarithm map of the
quotient wherever the minimiser is unique, which holds outside a set of measure zero. It is
the single-sample analogue of the minibatch optimal-transport couplings of
\citet{tong2024improving}, restricted to the permutation group. Three properties together
make the sampled distribution exchangeable: the aligned coupling, the exchangeable base
distribution (the prior), and a network that is equivariant to permutations of the source
tokens (\cref{app:architecture-equivariance}). The auxiliary endpoint loss uses the same
aligned logarithm, so it is invariant to permutations too (\cref{eq:aux-x}).

\subsection{Remarks}
\label{app:geometry-remarks}

\paragraph{Orientation.}
$(\psi,\iota,\phi_0)$ are Euler angles of the binary's frame relative to the line of sight, so
their natural space is $SO(3)$ modulo the symmetries of \cref{eq:psi-phi-symmetry}. Our chart
uses $\Sph^2\times\Sph^1$ instead, with $(\psi,\iota)$ as azimuth and polar angle. At the poles,
$\iota\in\{0,\pi\}$, the wave is circularly polarised and depends on $\psi$ and $\phi_0$ only
through $\phi_0\mp2\psi$, while the chart point no longer carries $\psi$. This set has
measure zero, but nearly face-on sources sit close to it, and an $SO(3)$ chart would treat
them more naturally.

\paragraph{Resolution of the coordinates in bfloat16.}
The network computes in bfloat16, which keeps 8 significant bits, so the input coordinate
$x_f$ is resolved to $2^{-8}|x_f|$ to $2^{-7}|x_f|$. One unit of $x_f$ spans half the window
interior, 24 bins on XS: the bfloat16 cell of $x_f$ ranges from below 0.01 bins near the
window centre to 0.094 bins near the edges. On XS this does not set the posterior width. At a
constant learning rate, the 0.088-bin floor did not depend on the cell size (Spearman
$-0.10$ over 379 loud sources). After the cooldown, the 768 network's loud sources in the
largest cells are $1.10\times$ (95\% CI 0.94--1.21) as wide as those in the smallest, and
the 512 network's $1.22\times$ (1.02--1.40). On B one unit spans 1536 bins, and the cell
reaches 6 bins over half of the window, about 160 times the XS floor; whether the network
still resolves sub-cell positions there is measured in \cref{app:b}.
\draftnote{New 2026-10-08: the XS cell test after the cooldown; not yet in the research log.
Decide whether it stays.}

% =======================================================================================
\section{Conditional Flow Matching from Prior to Posterior}
\label{app:fm}
% =======================================================================================

\subsection{Objective}
\label{app:fm-objective}

For a window $w$ with data $d$ and conditioning input $y = y(d, w)$ (\cref{app:conditioning}),
the target is the posterior $p(X\mid y, w)\propto p(d\mid X)\,\pi_w(X)$ over the set of sources
$X\in\manifold^S$ in flow coordinates. We learn a time-dependent vector field
$\vnet(X, t; y, w)$ whose flow transports the prior to the posterior: draws
$X(0)\sim\pi_w$, integrated along
\begin{equation}
  \frac{\dd X}{\dd t} = \dot s(t)\,\vnet\bigl(X, t; y, w\bigr),\qquad t\in[0,1],
  \label{eq:ode}
\end{equation}
should end at $X(1)\sim p(\cdot\mid y,w)$. The clock $s(t)$ is the identity in plain flow
matching; \cref{app:fm-warp} warps it.

Training uses the conditional flow-matching construction
\citep{lipman2023flow,liu2023flow,albergo2023building} on a Riemannian manifold
\citep{chen2024flow}. A training example draws a window $w$, a target set
$X_1\sim\pi_w$ and the data $d\sim p(d\mid X_1)$ from the simulator, and an independent base
set $X_0\sim\pi_w$ in the same window. The conditional path is the geodesic between the
aligned sets of \cref{eq:set-log},
\begin{equation}
  X_s = \Exp_{X_0}\bigl(s\,\Log_{X_0}(X_1)\bigr),\qquad
  u_s = \frac{\dd X_s}{\dd s},
  \label{eq:cond-path}
\end{equation}
and the network regresses the conditional velocity,
\begin{equation}
  \mathcal L_{\rm flow}(\eta) =
  \E_{t\sim\Unif(0,1)}\;\E_{w,\,X_0,\,X_1,\,d}\;
  \frac{1}{11S}\Bigl\lVert \vnet\bigl(X_{s(t)}, t; y, w\bigr) - u_{s(t)}\Bigr\rVert^2 .
  \label{eq:loss-flow}
\end{equation}
Its minimiser is the marginal field $v^\star(X,s;y,w) = \E[u_s\mid X_s = X, y, w]$, which
generates the marginal path from $p_0 = \pi_w$ (because $X_0$ is independent of the data) to
$p_1 = p(\cdot\mid y,w)$ \citep{lipman2023flow,chen2024flow}. The base distribution is thus the
prior itself, not a Gaussian: on most factors of $\manifold$ the prior is uniform
(\cref{app:geometry-chart}), and starting from it means that a source the data say nothing
about only needs to stay where it is. Flow-matching posterior estimation with a Gaussian base
is studied by \citet{dax2023flow}.

\subsection{A Warped Clock}
\label{app:fm-warp}

The fine structure of a narrow posterior is resolved only at the very end of the path
(\cref{app:sensitivity}), where a uniform clock spends almost none of its samples. We keep
$t\sim\Unif(0,1)$ and the network's time input $t$, but place the point at
\begin{equation}
  s(t) = 1 - (1-t)^p,\qquad p = 3,
  \label{eq:warp}
\end{equation}
along the path, and keep the target $u_s = \dd X_s/\dd s$, which stays of order one as
$s\to1$. Equivalently, the regression in \cref{eq:loss-flow} is weighted by the density of
path positions
\begin{equation}
  q_p(s) = \frac{1}{p}(1-s)^{1/p-1},\qquad
  \Pr\bigl(1-s<\epsilon\bigr) = \epsilon^{1/p}.
  \label{eq:warp-density}
\end{equation}
For $p=3$, 21\% of the samples fall in the first half of the path ($1-2^{-1/3}$, against 50\%),
and 1\% of them within $10^{-6}$ of its end (against $10^{-6}$). The minimiser is unchanged,
since only the weighting over $s$ moves, and the sampler restores the clock through the factor
$\dot s(t) = p(1-t)^{p-1}$ in \cref{eq:ode}. Conditioning the network on $t$ rather than $s$
stretches the end of the path: $1-s = 10^{-6}$ is $1-t = 10^{-2}$, which the sinusoidal time
embedding resolves (\cref{app:architecture}). This plays the role of the non-uniform noise
levels used to train diffusion models \citep{karras2022elucidating,kingma2023understanding,esser2024scaling}, applied here to the end of the path rather than to its middle.

\subsection{Auxiliary Heads and the Full Loss}
\label{app:fm-aux}

The network has two further heads (\cref{app:architecture}): an estimate $\hat X_1$ of the
endpoint and a reconstruction $\hat y$ of the noiseless conditioning input, with losses
\begin{equation}
  \mathcal L_{x} = \E\,\frac{1}{11S}\bigl\lVert\Log_{X_1}(\hat X_1)\bigr\rVert^2,
  \qquad
  \mathcal L_{y} = \E\,\frac{1}{\lvert y\rvert}\bigl\lVert\hat y - y_{\rm clean}\bigr\rVert^2,
  \label{eq:aux-x}
\end{equation}
where $\Log$ is the set-aligned logarithm of \cref{eq:set-log} and $y_{\rm clean}$ is the
conditioning input computed from the noiseless signal. The training loss is
\begin{equation}
  \mathcal L = \frac{\mathcal L_{\rm flow}}{V_{\rm flow}}
  + a(\tau)\left[\frac{\mathcal L_{x}}{V_{x}} + \frac{\mathcal L_{y}}{V_{y}}\right],
  \qquad
  a(\tau) = \tfrac12 + \tfrac12\cos\bigl(\pi\min(\tau/\tau_a, 1)\bigr),
  \label{eq:loss}
\end{equation}
with $V_\bullet$ the running variance of each target over all training samples so far,
pooled over coordinates (a Welford accumulator), and $a(\tau)$ a cosine decay over the first
fraction $f_{\rm aux} = \tau_a/\tau_{\rm total}$ of training ($\tau$ counts epochs of 1000
steps). The original runs used $f_{\rm aux} = 0.5$; the recipe uses $f_{\rm aux}=0.1$
(\cref{app:xs}). Once $a=0$ the auxiliary heads get no gradient, their losses drift
(\cref{app:xs-dynamics}), and sampling never uses them. The logged losses are the raw mean
squared errors before the variance normalisation.

\subsection{Sampling}
\label{app:fm-sampling}

To sample a window we draw $N$ base sets from the prior, integrate \cref{eq:ode} with the
classical fourth-order Runge--Kutta scheme on 32 uniform steps in $t$, and map back with
$\varphi_w^{-1}$ (\cref{alg:sample}). The four stages are evaluated in the ambient
coordinates and only the full step is retracted onto the manifold with $\Exp$. Under the
warped clock, steps uniform in $t$ crowd towards the end of the path: the last of 32 ends at
$1-s = 32^{-3} = 3\times10^{-5}$. Each draw costs $4\times32 = 128$ network evaluations. Draws
are independent, so we push them through in chunks of 256 to bound the memory on B. The
integrator is converged at 32 steps: 64 and 128 steps change neither the widths nor the
detections (\cref{tab:numerics}).

\begin{algorithm}[h]
  \caption{One training step (per example; batched over 256 windows).}
  \label{alg:train}
  \begin{algorithmic}
    \STATE {\bfseries Input:} weights $\eta$, auxiliary weight $a$, running variances $V_\bullet$
    \STATE draw a window $w$ (\cref{app:signal-windows}) and $\theta_{1:S}\sim\pi_w$
    \STATE $\tilde d \leftarrow \sum_s\tilde h(\theta_s) + \tilde n$ \hfill (\cref{eq:data})
    \STATE $y\leftarrow y(\tilde d, w)$,\ \ $y_{\rm clean}\leftarrow y(\tilde h, w)$ \hfill (\cref{app:conditioning})
    \STATE $X_1\leftarrow\varphi_w(\theta_{1:S})$,\ \ $X_0\leftarrow\varphi_w(\theta'_{1:S})$, $\theta'_{1:S}\sim\pi_w$
    \STATE $t\sim\Unif(0,1)$,\ \ $s\leftarrow1-(1-t)^p$
    \STATE $X_s\leftarrow\Exp_{X_0}\bigl(s\Log_{X_0}(X_1)\bigr)$,\ \ $u_s\leftarrow\dd X_s/\dd s$ \hfill (\cref{eq:set-log,eq:cond-path})
    \STATE $(\hat v, \hat X_1, \hat y)\leftarrow\text{network}_\eta(X_s, t, y, \fw)$
    \STATE update $V_\bullet$; take a gradient step on \cref{eq:loss} (Muon/Adam, \cref{app:training})
  \end{algorithmic}
\end{algorithm}

\begin{algorithm}[h]
  \caption{Sampling the posterior of one window.}
  \label{alg:sample}
  \begin{algorithmic}
    \STATE {\bfseries Input:} data $\tilde d$ of window $w$, number of draws $N$, steps $K_t=32$
    \STATE $y\leftarrow y(\tilde d, w)$;\ \ $\Delta t\leftarrow1/K_t$
    \FOR{$n=1$ {\bfseries to} $N$ (in parallel)}
      \STATE $X\leftarrow\varphi_w(\theta_{1:S})$, $\theta_{1:S}\sim\pi_w$
      \FOR{$j=0$ {\bfseries to} $K_t-1$}
        \STATE $t\leftarrow j\Delta t$;\ \ $F(X,t) := p(1-t)^{p-1}\,\vnet(X,t;y,\fw)$
        \STATE $k_1\leftarrow F(X,t)$,\ \ $k_2\leftarrow F(X+\tfrac{\Delta t}{2}k_1, t+\tfrac{\Delta t}{2})$
        \STATE $k_3\leftarrow F(X+\tfrac{\Delta t}{2}k_2, t+\tfrac{\Delta t}{2})$,\ \ $k_4\leftarrow F(X+\Delta t\,k_3, t+\Delta t)$
        \STATE $X\leftarrow\Exp_X\bigl(\tfrac{\Delta t}{6}(k_1+2k_2+2k_3+k_4)\bigr)$
      \ENDFOR
      \STATE {\bfseries output} $\varphi_w^{-1}(X)$
    \ENDFOR
  \end{algorithmic}
\end{algorithm}

% =======================================================================================
\section{Why the Flow-Matching Loss Cannot See a Narrow Posterior}
\label{app:sensitivity}
% =======================================================================================

The first trained XS models placed loud sources at the right frequency, but with an $f_0$
width that stopped shrinking at about $0.15$--$0.2$ bins whatever the SNR: about 100 times the
ideal width at SNR 1000 (\cref{fig:width-snr}). The training loss gave no sign of it. This
appendix explains why, with a model simple enough to solve, and derives what the warped clock
of \cref{eq:warp} changes.

\subsection{A Gaussian Model of One Coordinate}

Take one flow coordinate with base $x_0\sim\Normal(0,s_0^2)$, target posterior
$x_1\sim\Normal(0,\sigma^2)$, the straight path $x_s = (1-s)x_0 + s x_1$ and the target
$u = x_1 - x_0$. At path position $s$ the marginal is $x_s\sim\Normal(0, m^2(s))$ with
$m^2 = (1-s)^2s_0^2 + s^2\sigma^2$, and the optimal field is linear,
\begin{equation}
  v^\star(x,s) = \E[u\mid x_s = x] = g(s)\,x,\qquad
  g(s) = \frac{s\,\sigma^2 - (1-s)\,s_0^2}{m^2(s)} .
\end{equation}
A model whose field is exactly that of a wider posterior $\Normal(0,\sigma_m^2)$, the field
$g_m(s)\,x$ with $\sigma\to\sigma_m$ in $g$ and $m$, pays an excess loss over the optimum. By
the orthogonality of the regression residual, it is
\begin{equation}
  \Delta L_p(\sigma_m;\sigma) = \E_{t}\,\E_{x\sim\Normal(0,m^2)}\bigl[(g_m - g)^2x^2\bigr]
  = \int_0^1 m^2(s)\,\bigl[g_m(s) - g(s)\bigr]^2\,\dd t,\qquad s = 1-(1-t)^p .
  \label{eq:excess}
\end{equation}

\begin{proposition}
\label{prop:sensitivity}
For $\sigma\ll\sigma_m\ll s_0$,
\begin{equation}
  \Delta L_p(\sigma_m;\sigma) = C_p\,s_0^{\,2-1/p}\,\sigma_m^{1/p}\,\bigl(1+o(1)\bigr),\qquad
  C_p = \frac{(2p-1)\,\pi}{4p^2\sin\bigl(\pi/2p\bigr)},
  \label{eq:sensitivity}
\end{equation}
so that $C_1 = \pi/4$ for the uniform clock and $C_3 = 5\pi/18$ for the warped one. To leading
order the excess does not depend on the true width $\sigma$.
\end{proposition}

\begin{proof}
Set $s_0=1$ and $\epsilon = 1-s$; the clock has density $q_p(\epsilon) = \frac1p\epsilon^{1/p-1}$
in $\epsilon$. Where $\epsilon\gg\sigma$, $m^2\simeq\epsilon^2$ and $g\simeq-1/\epsilon$, so
\begin{equation*}
  g_m - g \simeq \frac{\sigma_m^2-\epsilon}{\epsilon^2+\sigma_m^2} + \frac1\epsilon
  = \frac{\sigma_m^2(1+\epsilon)}{\epsilon\,(\epsilon^2+\sigma_m^2)},\qquad
  m^2(g_m-g)^2 \simeq \frac{\sigma_m^4}{(\epsilon^2+\sigma_m^2)^2}
\end{equation*}
up to a relative $O(\epsilon)$. The region $\epsilon\lesssim\sigma$ contributes $O(\sigma)$.
Substituting $\epsilon = \sigma_m z$,
\begin{equation*}
  \Delta L_p \simeq \frac1p\int_0^\infty \epsilon^{1/p-1}\frac{\sigma_m^4\,\dd\epsilon}{(\epsilon^2+\sigma_m^2)^2}
  = \frac{\sigma_m^{1/p}}{p}\int_0^\infty\frac{z^{1/p-1}}{(1+z^2)^2}\,\dd z
  = \frac{\sigma_m^{1/p}}{2p}\,B\!\Bigl(\frac{1}{2p}, 2-\frac{1}{2p}\Bigr),
\end{equation*}
and $\Gamma(a)\Gamma(2-a) = (1-a)\,\pi/\sin(\pi a)$ with $a = 1/2p$ gives $C_p$. The factor
$s_0^{2-1/p}$ follows from dimensions: $\Delta L$ scales as $s_0^2$ and depends on
$\sigma_m/s_0$ alone.
\end{proof}

\Cref{fig:warp}b checks \cref{eq:sensitivity} against direct quadrature of \cref{eq:excess}:
the ratio is 1.000 for $p=1$ and 0.95 for $p=3$ at a width of 0.037 XS bins.

\begin{figure}[h]
  \centering
  \includegraphics[width=0.55\textwidth]{figures/warp.pdf}
  \caption{(a) Share of the training samples that land within $\epsilon$ of the end of the
  path, $\epsilon^{1/p}$, for the uniform and the warped clock; vertical lines mark where the
  structure of an $f_0$ posterior of 0.18 bins, 0.037 bins, and the ideal width at SNR 1000
  is resolved ($1-s\approx\sigma/s_0$, XS units). (b) Excess flow-matching loss of a model that
  is $\sigma_m$ wide when the true posterior is far narrower: quadrature of \cref{eq:excess}
  (dots) and the asymptote \cref{eq:sensitivity} (lines).}
  \label{fig:warp}
\end{figure}

\subsection{Consequences}

\paragraph{The loss rewards absolute narrowness, weakly.}
Under the uniform clock the excess is linear in the model's width,
$\Delta L_1\simeq\frac\pi4 s_0\sigma_m$, and its gradient $\frac\pi4 s_0$ is the same whether
the model is 1000 or 2 times too wide. Nothing in the objective grows as the posterior gets
closer to its true width.

\paragraph{The floor is invisible in the logged loss.}
On XS, $x_f$ is uniform on $[-1,1]$, so $s_0 = 1/\sqrt3$, and one unit is 24 bins. A floor of
0.18 bins, $\sigma_m = 0.0075$, costs $\frac\pi4 s_0\sigma_m = 3.4\times10^{-3}$ on that
coordinate. The logged loss averages over $11S = 44$ coordinates, and about half of the
sources are loud enough to be affected, so the whole floor raises the logged loss by
$\approx1.5\times10^{-4}$. That is less than the epoch-to-epoch scatter of the logged loss
($5\times10^{-4}$), and about 1\% of what a doubling of the training steps buys
($1.7\times10^{-2}$, \cref{app:xs-dynamics}). Longer training and bigger networks lower the
part of the loss the objective can see, which is the coarse structure: detection, sky
position, amplitude.

\paragraph{Where the information lives.}
The structure of a posterior of width $\sigma$ is resolved only at $1-s\lesssim\sigma/s_0$. For
the 0.18-bin floor this is the last 1.3\% of a uniform clock, and for the ideal width at SNR
1000 ($5.5\times10^{-4}$ bins) the last $4\times10^{-5}$ (\cref{fig:warp}a).

\paragraph{What the warp changes.}
With $p=3$ the same floor costs $(10/9)(\sigma_m/s_0)^{-2/3}$ times more: 20 times at 0.18
bins and 58 times at 0.037 bins. The gradient now grows as the posterior narrows,
$\partial_{\sigma_m}\Delta L_3\propto\sigma_m^{-2/3}$, and 23\% of the samples land where the
0.18-bin structure is resolved instead of 1.3\%. In the experiments the warp alone narrowed the
loud-source $f_0$ posteriors by a factor of 0.6 at every SNR (\cref{app:xs}).

The model is idealised: one coordinate, Gaussian base and posterior, a straight path, a model
error of the simplest form. Its point is the scaling. A uniform-$t$ flow-matching loss is
nearly blind to the width of posteriors much narrower than the prior, which is the regime of
loud gravitational-wave sources, whose ideal posterior is three to seven orders of magnitude
narrower than the window (48 bins on XS and 3072 on B, against $5.5\times10^{-4}$ bins at SNR
1000).

% =======================================================================================
\section{The Conditioning Input}
\label{app:conditioning}
% =======================================================================================

The network sees the data of a window through three operations: whitening, a
Wilson--Daubechies--Meyer (WDM) transform, and a compression of the dynamic range.

\paragraph{Whitening.}
Each bin and channel is divided by its noise amplitude,
$\tilde d^{\,\rm w}_{k,c} = \tilde d_{k,c}/\sqrt{S_c(f_k)\Tobs/2}$, so that under noise alone
$\tilde d^{\,\rm w}_{k,c}\sim\CN(0,1)$: real and imaginary parts have variance $1/2$, and the
signal enters in units of its SNR per bin.

\paragraph{WDM transform.}
The WDM transform \citep{necula2012transient,cornish2020time,vajpeyi2026explicit} tiles the
time--frequency plane with $N_t$ time rows and $N_f$ frequency channels. From the one-sided
spectrum, channel $m$ draws on the Fourier bins within $N_t/2$ of $mN_t/2$:
\begin{equation}
  x_{n,m} = \sum_{l=-N_t/2}^{N_t/2-1}\tilde d^{\,\rm w}_{mN_t/2+l}\,\Phi_l\,e^{2\pi\ii nl/N_t},
  \qquad
  w_{n,m} = \sqrt2\,(-1)^{nm}\,\mathrm{Re}\bigl[C^*_{n,m}\,x_{n,m}\bigr],
  \qquad
  C_{n,m} = \begin{cases}1 & n+m\ \text{even}\\ \ii & n+m\ \text{odd},\end{cases}
  \label{eq:wdm}
\end{equation}
with $\Phi$ a Meyer-type window with a flat top over $|2l/N_t|\le a$ and a cosine taper of
width $b = 1-2a$, here $a = 1/3$. We compute only the channels that cover the window: $K$
channels starting at an even offset, as one reshape and a small matrix product, so that the
compiled graph does not grow with $K$ (a channel-by-channel loop compiles to millions of
operations on B).

We use the minimum $N_t = 2$ time rows, and in this case the transform simplifies exactly. The
window evaluates to $\Phi = (1, 0)$, so channel $m$ sees only bin $m$, and \cref{eq:wdm}
reduces to
\begin{equation}
  (w_{0,m},\,w_{1,m}) =
  \begin{cases}
    \sqrt2\,\bigl(\mathrm{Re}\,\tilde d^{\,\rm w}_m,\ \ \mathrm{Im}\,\tilde d^{\,\rm w}_m\bigr) & m\ \text{even},\\
    \sqrt2\,\bigl(\mathrm{Im}\,\tilde d^{\,\rm w}_m,\ -\mathrm{Re}\,\tilde d^{\,\rm w}_m\bigr) & m\ \text{odd},
  \end{cases}
  \label{eq:wdm-nt2}
\end{equation}
which we verified numerically against the implementation. The conditioning image is therefore
a lossless, real re-encoding of the whitened complex spectrum, two rows of unit-variance
pixels per TDI channel, rather than a time--frequency map. More time rows would trade
frequency resolution for time resolution, which the slowly evolving signals of this problem
do not need inside a window.
\draftnote{With $N_t=2$ the WDM step is a re-encoding of Re/Im. Present it that way in the
main text, or justify the WDM framing (e.g.\ a path to longer windows or other sources).}

\paragraph{Compression.}
A source puts $\sum_{k,c}|\tilde h^{\rm w}_{k,c}|^2 = \snr^2/2$ into a few bins, so the loudest
reach whitened amplitudes of order $10^3$ per bin, against unit-variance noise.
The input is $y = \operatorname{arcsinh}(w)$, linear in the noise and logarithmic for loud
signals, arranged as an image of $2\times K\times3$ pixels: time rows, frequency channels, and
the $A$, $E$, $T$ channels (\cref{fig:window}b). The same operations applied to the noiseless
signal give the target $y_{\rm clean}$ of the reconstruction head (\cref{eq:aux-x}).

\paragraph{What the image does not carry.}
Whitening and windowing remove the absolute frequency, which still matters: the Doppler width
scales with $f$, the chirp with $f_0^{11/3}$, and the amplitude window with $S_n(\fw)$. The
window centre $\fw$ is therefore given to the network separately (\cref{app:architecture}).

\begin{figure}[h]
  \centering
  \includegraphics[width=\textwidth]{figures/window.pdf}
  \caption{One XS window, from data to posterior (eval injection q0.20, final model). (a)
  Whitened power per bin, summed over $A$ and $E$, with the noiseless signal; vertical lines
  mark the true frequencies and their SNRs, and the shaded bins are the guard bands. (b) The
  conditioning input: $\operatorname{arcsinh}$ of the two-row WDM coefficients
  (\cref{eq:wdm-nt2}) for the three TDI channels. (c) The flow's $f_0$ posterior for each
  source, relabelled onto it (\cref{app:evaluation-matching}), against the Fisher forecast
  (each histogram scaled to a unit peak). The three detectable sources are found, with
  posteriors several times wider than the Fisher forecast. The SNR-5.3 source is not found,
  and its posterior does not spread over the window as the prior would: with exactly four
  sources per window and amplitudes of SNR $\gtrsim7$, a fourth source the data do not support
  can only sit where it is masked by the loud ones, which is where the flow puts it.}
  \label{fig:window}
\end{figure}

% =======================================================================================
\section{Network Architecture}
\label{app:architecture}
% =======================================================================================

The velocity field is a transformer with two token streams, one for the sources and one for
the data, that attend to each other in every block: the multimodal diffusion transformer
(MMDiT) of \citet{esser2024scaling}, with the adaptive-normalisation conditioning of
\citet{peebles2023scalable}. \Cref{fig:architecture} shows it, and \cref{tab:network} gives
the sizes.

\begin{figure}[h]
  \centering
  \includegraphics[width=\textwidth]{figures/architecture.pdf}
  \caption{The network. The parameter stream (blue) has one token per source and no
  positional embedding, so the network is equivariant to permutations of the sources; the
  data stream (grey) has one token per $2\times2$ patch of the conditioning image. Both
  streams attend jointly in every block, and the flow time and the window frequency modulate
  every normalisation and residual gate. Only the velocity head is used to sample; the two
  auxiliary heads are trained during the first 10\% of the steps.}
  \label{fig:architecture}
\end{figure}

\subsection{Embeddings}

All feed-forward maps are SwiGLU \citep{shazeer2020glu},
$\mathrm{FF}(z) = W_2\bigl(a\odot\mathrm{SiLU}(g)\bigr)$ with $[a;g] = W_1z + b_1$.
\begin{itemize}
  \item \textbf{Sources.} Each source's 11 coordinates are embedded by its own application of
        a shared $\mathrm{FF}: \R^{11}\to\R^D$, giving $S$ tokens with no positional embedding.
  \item \textbf{Data.} The $2\times K\times3$ image is cut into $2\times2$ patches, each
        flattened to 12 numbers and mapped linearly to $\R^D$, giving $2R$ tokens (1 row of
        $K/2$). Sinusoidal embeddings of the row and column index, each passed through an
        $\mathrm{FF}$, are added.
  \item \textbf{Conditioning.} The flow time and the window centre frequency are embedded as
        $c = \mathrm{FF}_t\bigl(E(t)\bigr) + \mathrm{FF}_f\bigl(E(\fw)\bigr)$, with
        $E(\tau) = \mathrm{FF}\bigl([\sin2\pi\omega_k\tau;\ \cos2\pi\omega_k\tau]_{k=1}^{D}\bigr)$
        and $\omega_k = (2\pi)^{-(k-1)/(D-1)}$, angular frequencies from 1 to $2\pi$ per
        unit. For $\fw$ in hertz the phases are small and the embedding is close to linear in
        $\fw$.
\end{itemize}

\subsection{Blocks}

Each of the $N$ blocks updates the source tokens $z^x$ and the data tokens $z^y$. For each
stream $\sigma\in\{x,y\}$, a zero-initialised linear map of $\mathrm{SiLU}(c)$ produces a shift,
a scale and a gate, $(\beta^\sigma,\gamma^\sigma,\zeta^\sigma)$, so that every residual branch
starts as the identity \citep{peebles2023scalable}:
\begin{align}
  h^\sigma &= \mathrm{LN}(z^\sigma)\odot(1+\gamma^\sigma) + \beta^\sigma, \nonumber\\
  [q^\sigma; k^\sigma; v^\sigma] &= W^\sigma_{qkv}h^\sigma,\qquad
  o = \mathrm{Attn}\bigl([q^x; q^y], [k^x; k^y], [v^x; v^y]\bigr), \nonumber\\
  z^\sigma &\leftarrow z^\sigma + \zeta^\sigma\odot W^\sigma_{o}\,o^\sigma, \nonumber\\
  z^\sigma &\leftarrow z^\sigma + \zeta'^\sigma\odot\mathrm{FF}^\sigma\bigl(\mathrm{LN}(z^\sigma)\odot(1+\gamma'^\sigma) + \beta'^\sigma\bigr),
\end{align}
with a second shift, scale and gate $(\beta'^\sigma,\gamma'^\sigma,\zeta'^\sigma)$ for the
feed-forward branch, $\mathrm{LN}$ a layer normalisation without learned affine parameters,
and the feed-forward hidden width $2D$. The attention runs over all $S+2R$ tokens with $H$
heads of dimension 64. Queries and keys are RMS-normalised per head, with no learned scale
\citep{henry2020query,dehghani2023scaling}. The projections are separate for the two streams,
the attention is joint. The blocks are applied with a scan over stacked parameters and
rematerialised in the backward pass \citep{chen2016training}.

\subsection{Heads}

The source tokens pass through a final shift-and-scale modulation and an
$\mathrm{FF}:\R^D\to\R^{2\times11}$, which gives the velocity $\vnet$ and the endpoint
estimate $\hat X_1$ for each source. The data tokens are mapped back to $2\times2$ patches,
which gives the reconstruction $\hat y$.

\subsection{Permutation Equivariance}
\label{app:architecture-equivariance}

Every map applied to a source token is shared across sources, and attention without
positional information is equivariant to permutations of its tokens. Permuting the input
sources therefore permutes the outputs: $\vnet(\sigma X) = \sigma\,\vnet(X)$ for every
$\sigma\in\mathfrak S_S$. With an exchangeable base distribution, the flow of \cref{eq:ode}
maps exchangeable distributions to exchangeable distributions, and the set alignment of
\cref{eq:set-log} gives it a regression target that is itself defined on unordered sets
\citep[compare][]{zaheer2017deep,lee2019set}.

\subsection{Sizes and Precision}

\Cref{tab:network} gives the two widths we trained. The network is the same on every rung,
and only the number of data tokens changes. Parameters are kept in float32, and the forward
and backward passes run in bfloat16 \citep{micikevicius2018mixed,kalamkar2019study}. In
bfloat16, attention uses the fused cuDNN flash-attention kernel \citep{dao2022flashattention}:
on B, materialising the attention logits of a batch, $256\times8$ heads $\times\,2052^2$ in
float32, would need 32\,GiB, which is what made the first B run run out of memory.

\begin{table}[h]
\caption{Network sizes. Parameter counts are exact; tokens are sources plus data patches; step times are for a batch of 256 windows (\cref{app:training-compute}).}
\label{tab:network}
\centering
\small
\begin{tabular}{lrrrrr}
\toprule
width $D$ & heads $H$ & blocks $N$ & parameters & tokens on XS / S / B & A100 time per step, XS / B \\
\midrule
512 & 8 & 8 & 75\,636\,258 & 36 / 132 / 2052 & 38.9\,ms / 1.54\,s \\
768 & 12 & 8 & 170\,077\,474 & 36 / 132 / 2052 & 59.6\,ms / 2.57\,s \\
\bottomrule
\end{tabular}
\end{table}

% =======================================================================================
\section{Optimisation, Schedules and Compute}
\label{app:training}
% =======================================================================================

\subsection{Optimiser}

Gradients are clipped to a global norm of 1, then the two-dimensional weight matrices are
updated with Muon \citep{jordan2024muon,liu2025muon} and all other parameters with Adam, both
at learning rate $\gamma = 10^{-4}$ and without weight decay (\texttt{optax.contrib.muon},
version 0.2.8, default width scaling). Batches hold 256 windows; one epoch is 1000 steps.
Muon orthogonalises each matrix update, so its size per element is fixed by the learning rate
and the shape: we measured an update RMS of $(0.03$--$0.045)\,\gamma$, about five times smaller
than AdamW at the same nominal rate. We did not tune the learning rate; with the compute
available, every run went into a different question (\cref{app:xs}).

\subsection{Schedules}
\label{app:training-schedules}

\paragraph{Auxiliary heads.}
The weight $a(\tau)$ of \cref{eq:loss} decays as a half cosine from 1 to 0 over the first
$f_{\rm aux}$ of training: half of the run in the original configuration, a tenth in the
recipe.

\paragraph{Learning rate.}
The original runs kept the learning rate constant. The recipe adds a cooldown: constant, then
a linear decay to zero over the last 20\% of the steps, the warmup-stable-decay schedule of
\citet{hu2024minicpm} and \citet{hagele2024scaling} without warmup. It is applied as a scale
on the update in epoch $e$ of $E$,
\begin{equation}
  \lambda_e = \begin{cases} 1 & e < E - E_c,\\[2pt] \dfrac{E - e - \frac12}{E_c} & e\ge E-E_c,\end{cases}
\end{equation}
with $E_c$ the number of cooldown epochs, so that each epoch takes the ramp's value at its
midpoint. Without weight decay, scaling the update is exactly scaling the learning rate, and
because the scale is a function of the epoch rather than of a step counter in the optimiser
state, a run resumed with a fresh optimiser lands at the right point of the schedule. The
cooldowns of \cref{app:xs} continue a finished 1M-step run for 200k more steps with
$E_c = 200$.

\paragraph{Checkpoints.}
Every epoch saves a checkpoint, and a job longer than the cluster's 24-hour limit is a chain
of jobs that each resume from the latest one. To fit a storage quota, the 768-wide runs and
B save the weights without the optimiser state. A resumed job then restarts Muon's momentum
and Adam's moments, which leaves no visible trace in the loss: on B, the epoch before and
after the first restart logged 0.5370 and 0.5374.

\subsection{Compute}
\label{app:training-compute}

Every run used a single NVIDIA A100 (80\,GB, PCIe). The simulator, the loss and the update
are compiled into one XLA program per epoch of 1000 steps. \Cref{tab:compute} gives the
measured costs. A fit of the step time against the network's FLOPs on XS and S gives about
23\,ms of fixed cost per step and 105\,TFLOP/s of marginal throughput (about a third of the
A100's bfloat16 peak), so the small XS problem is dominated by overheads and the large B
problem by the network. The whole XS ablation ladder of \cref{tab:xs-ladder} cost about 45
A100-hours, and B1 about 86.

\begin{table}[h]
\caption{Measured cost on one A100, batch 256. Peak memory is for training; JAX reserves 75\%
of the card by default.}
\label{tab:compute}
\centering
\small
\begin{tabular}{llrrl}
\toprule
rung & width & per 1000 steps & peak memory & runs \\
\midrule
XS & 512 & 38.9\,s & -- & 1M steps in 10.8\,h; 200k cooldown in 2.2\,h \\
XS & 768 & 59.6\,s & -- & 1M steps in 16.6\,h; 200k cooldown in 3.6\,h \\
S & 512 & 80\,s & -- & 500k steps in 11.1\,h (an early run that diverged; not retrained) \\
B & 512 & 1536\,s & 18.6\,GiB & 200k steps in 86\,h (4 chained jobs) \\
B & 768 & 2570\,s & 29.4\,GiB & 300-step benchmark \\
\bottomrule
\end{tabular}
\end{table}

\paragraph{Sampling cost.}
A scorecard pass on XS (210 windows, 256 draws each, 32 RK4 steps) takes about 4 minutes on
an A100, about one second per window. On B, 256 draws of one window take about 78 seconds with
the 768-wide network.

\subsection{Numerics}
\label{app:training-numerics}

Training and evaluation run in bfloat16. Three checks, on XS-late
(\cref{tab:numerics}), show that the numerics do not set the results: the RK4 integrator is
converged at 32 steps; float32 evaluation on a CPU gives the same widths as bfloat16 on the
A100 (median ratio 0.991, 95\% interval 0.978--1.000, over 93 loud sources) and finds the same
169 sources; and float32 on the A100 reproduces the CPU once XLA's GEMM autotuning is switched
off. With autotuning on, float32 programs were silently wrong on both GPUs we used. At the
default level they found almost no sources. At level 3 the A100 gave posteriors 2.4 times
narrower than the CPU's at the same distance from the truth: over-confident, with 45\% of the
truths inside the 68\% intervals. A fixed compiled program gave different answers at different
autotuning levels, on two JAX and two CUDA versions, so the fault is in the compilation and not
in the model; bfloat16 is unaffected. We report it because the over-confident failure is easy to
mistake for a sharper posterior.

% =======================================================================================
\section{Evaluation Protocol and Metrics}
\label{app:evaluation}
% =======================================================================================

\subsection{Injections}

All models are scored on the same test set. With a fixed seed we draw 1024 candidate windows
and their sources from the prior. The ten \emph{eval injections} are the candidates at the
0.1, 0.2, \dots, 1.0 quantiles of the window SNR (121 to 1848 on XS), and 200 further
candidates are taken at random. This gives 210 windows and 840 sources. Each window gets its
own noise realisation, and the flow draws 256 samples per window (1024 for corner plots) with
32 RK4 steps in bfloat16. Every model sees the same windows, noise and base draws, so
comparisons between models are paired, source by source. These are prior-predictive
injections: averaged over them, a calibrated posterior gives uniform ranks
\citep{talts2018validating}.

\subsection{Relabelling the Draws}
\label{app:evaluation-matching}

The flow's source slots are exchangeable, so the first slot of one draw can be any of the true
sources. To score a source, each draw is relabelled by the permutation that minimises
$\sum_i\lVert F(\hat\theta_{\sigma(i)}) - F(\theta_i)\rVert^2$, with features
$F = \bigl(f_0/(f_{\rm hi}-f_{\rm lo}),\ \log_{10}\Amp,\ n_{\rm sky}\bigr)$. The matching weighs
sky position and amplitude as well as frequency. When a window also holds a source the flow
has not localised, the matching can give a well-localised source's slot to that one: in the
window of \cref{fig:corner-median}, 18\% of the draws are relabelled this way, although 99.7\%
of them have a slot within one bin of the source. The robust width below is insensitive to this
contamination, while the coverage it affects is, if anything, made conservative.

\subsection{Metrics}
\label{app:evaluation-metrics}

All metrics are computed on the $f_0$ marginal of each source, in bins of $1/\Tobs$:
\begin{itemize}
  \item \textbf{width}: $w = 1.4826\,\MAD$ of the draws, the standard deviation of a Gaussian,
        robust to a minority of mis-relabelled draws;
  \item \textbf{found}: $w<1$ bin. A source the flow has not found is spread over the window,
        several bins wide;
  \item \textbf{ideal width}: $w^\star = \sqrt3/(\pi\snr)$, with $\snr$ the source's own SNR
        (\cref{eq:ideal-width}). On the 40 eval sources, the full Fisher forecast below is a
        median 10\% wider (0 to 33\%), so ``$\times$ ideal'' reads about 10\% higher than
        ``$\times$ Fisher'';
  \item \textbf{offset}: median of the draws minus the truth;
  \item \textbf{rank}: the fraction of draws below the truth, uniform on $[0,1]$ for a
        calibrated posterior; a mean away from 1/2 is a bias;
  \item \textbf{coverage}: whether the truth lies in the central 68\% and 95\% intervals
        (in68, in95).
\end{itemize}
We report them in SNR bands ($<15$, 15--40, 40--100, $\ge100$). Widths and offsets are medians
over the found sources of a band, and coverage and ranks are averages over all its sources.

\paragraph{The ideal width.}
For a sinusoid $s(t) = A\cos(\omega t+\phi)$ observed for a time $T$ in white noise, and a
time origin at the centre of the observation, the Fisher matrix in $(\omega,\phi)$ is
diagonal, with $\Gamma_{\omega\omega} = \snr^2T^2/12$ and $\Gamma_{\phi\phi} = \snr^2$. The
amplitude decouples. The frequency error is therefore
\begin{equation}
  \sigma_{f_0} = \frac{\sqrt{12}}{2\pi\,\snr\,T} = \frac{\sqrt3}{\pi\,\snr}\ \text{bins},
  \label{eq:ideal-width}
\end{equation}
0.055 bins at SNR 10 and $5.5\times10^{-4}$ bins at SNR 1000. Orbital Doppler modulation,
the chirp and correlations with the other seven parameters can only widen it: a lower bound,
attained by the full Fisher forecast to within the 10\% above.

\subsection{Fisher Forecast}
\label{app:evaluation-fisher}

The corner plots overlay a Gaussian forecast at the truth. Its precision is the Hessian of
the negative log-likelihood (\cref{eq:loglik}) at the truth, averaged over 32 noise
realisations, plus the prior precision, taken from the covariance of the 1024 candidates. A
single realisation's Hessian can be indefinite along poorly constrained directions; the
average cancels the data-dependent term. In physical units the precision matrix spans about 45
orders of magnitude on its diagonal, so we diagonalise it after a Jacobi rescaling to unit
diagonal, floor each eigenvalue at the prior precision along its eigenvector, and draw from the
resulting Gaussian. The forecast is a reference width, not a posterior
\citep{cutler1994gravitational,vallisneri2008use}. It is unimodal, so it cannot show the
exact symmetries of \cref{eq:psi-phi-symmetry}, and Gaussian, so it cannot follow the curved
amplitude--inclination degeneracy. For a parameter the data do not constrain, such as the
chirp mass at low frequency, where the chirp drift over the mission is a small fraction of a
bin (0.08 bins at 2\,mHz for $\Mc = 0.45\,\Msun$), it simply shows the prior's width centred on
the truth.

% =======================================================================================
\section{Experiments on XS}
\label{app:xs}
% =======================================================================================

\subsection{Setup}

XS is the smallest rung: $R=16$, 64-bin windows with a 48-bin interior, $f_0$ between 0.1 and
1.26\,mHz (53\% of the log band), four sources per window and 32 data tokens
(\cref{tab:rungs}). Because a step is cheap (39\,ms at width 512), XS is where every design
decision was made. With a tight compute budget, we never ran single-variable sweeps. Each
run combined the changes we believed in, read against the runs already paid for, and was
followed by cheap evaluation-only tests that ruled out alternative explanations. The five
models of \cref{tab:xs-ladder} form the ladder of this section:
\begin{enumerate}
  \item \textbf{uniform clock}: width 512, uniform clock ($p=1$), auxiliary heads over the
        first half of training ($f_{\rm aux}=0.5$), constant learning rate, 500k steps then
        continued to 1M;
  \item \textbf{warp}: the warped clock ($p=3$) and a short auxiliary phase
        ($f_{\rm aux}=0.1$), otherwise identical, 1M steps;
  \item \textbf{warp, 768}: the same at width 768 (170M parameters);
  \item \textbf{warp, cooldown}: model 2 continued for 200k steps with the learning rate
        decaying linearly to zero;
  \item \textbf{warp, 768, cooldown} (\emph{final}): model 3 continued the same way.
\end{enumerate}
The evaluation is that of \cref{app:evaluation}: 840 sources in 210 windows, scored
identically for every model.

\begin{table}[h]
\caption{The XS ladder, on the same 840 sources. Loud width: median $f_0$ width of the found
sources at SNR $\ge100$, and its median ratio to the ideal width $\sqrt3/(\pi\snr)$. Found:
fraction of sources with width below one bin. Calibration: coverage of the 68\% and 95\%
intervals and mean rank of the truth at SNR $\ge100$. Bias: the largest median offset across
eight slices of the source's position in the window, SNR $\ge40$ (\cref{fig:bias}). Hours: A100
time, cooldowns added to the run they continue.}
\label{tab:xs-ladder}
\centering
\small
\setlength{\tabcolsep}{4pt}
% generated by paper/scripts/tables.py from the scorecards; do not edit by hand
\begin{tabular}{llllcccccccccl}
\toprule
& & & & \multicolumn{2}{c}{loud $f_0$ width ($\rho\geq100$)} & \multicolumn{3}{c}{found fraction, $\rho$:} & \multicolumn{3}{c}{calibration, $\rho\geq 100$} & bias & A100 \\
\cmidrule(lr){5-6} \cmidrule(lr){7-9} \cmidrule(lr){10-12}
clock & width & steps & lr & bins & $\times$ ideal & $<15$ & 15--40 & 40--100 & in68 & in95 & rank & [bins] & hours \\
\midrule
uniform & 512 & 1M & const. & 0.147 & 94 & 0.52 & 0.85 & 0.95 & 0.77 & 0.97 & 0.62 & 0.146 & 10.8 \\
warp & 512 & 1M & const. & 0.088 & 59 & 0.45 & 0.81 & 0.92 & 0.87 & 0.99 & 0.50 & 0.045 & 10.9 \\
warp & 768 & 1M & const. & 0.087 & 55 & 0.47 & 0.83 & 0.94 & 0.82 & 0.98 & 0.41 & 0.081 & 16.6 \\
warp & 512 & 1.2M & cooled & 0.039 & 26 & 0.50 & 0.85 & 0.95 & 0.93 & 0.99 & 0.47 & 0.009 & 10.9 + 2.2 \\
warp & 768 & 1.2M & cooled & 0.037 & 25 & 0.54 & 0.86 & 0.97 & 0.93 & 0.99 & 0.49 & 0.008 & 16.6 + 3.6 \\
\bottomrule
\end{tabular}

\end{table}

\subsection{Headline Results}

The final model finds 54\%, 86\%, 97\% and 99.5\% of the sources in the four SNR bands
(\cref{tab:xs-final}), places them without bias (median offsets below 0.006 bins, mean ranks
0.49--0.53), and is calibrated on the conservative side: 78--93\% of the truths inside the
68\% intervals and 97--100\% inside the 95\% intervals. Below SNR 40 its $f_0$ posteriors are
within a factor 2--3 of the ideal width. Above that, the width stops following $1/\snr$
(\cref{fig:width-snr}): it levels off at $0.035$--$0.039$ bins from SNR 100 to 1400, 25 times
the ideal width at the median loud source. This floor, 0.6\,nHz, is the main remaining
limitation, and the rest of this section traces what moved it.

\begin{table}[h]
\caption{The final model (XS, width 768, warped clock, cooldown) by SNR band, on 840 sources.
Width and offset are medians over found sources, in bins of $1/\Tobs$; coverage and rank
over all sources.}
\label{tab:xs-final}
\centering
\small
% generated by paper/scripts/tables.py from the scorecards; do not edit by hand
\begin{tabular}{lrcccrccc}
\toprule
SNR $\rho$ & $n$ & found & width [bins] & $\times$ ideal & offset [bins] & in68 & in95 & rank \\
\midrule
$<15$ & 148 & 0.54 & 0.113 & 2 & -0.006 & 0.78 & 0.97 & 0.53 \\
15--40 & 157 & 0.86 & 0.068 & 3 & -0.003 & 0.90 & 1.00 & 0.50 \\
40--100 & 151 & 0.97 & 0.041 & 5 & +0.000 & 0.90 & 0.99 & 0.49 \\
$\geq 100$ & 384 & 0.99 & 0.037 & 25 & +0.002 & 0.93 & 0.99 & 0.49 \\
\midrule
all & 840 & 0.89 & & & & 0.89 & 0.99 & 0.50 \\
\bottomrule
\end{tabular}

\end{table}

\begin{figure}[h]
  \centering
  \includegraphics[width=0.6\textwidth]{figures/width_snr.pdf}
  \caption{Width of the $f_0$ posterior against the source's SNR, on the 840 scored sources.
  Points: the final model. Lines: medians of the found sources in SNR bins, for the
  uniform-clock baseline, the warped clock, and the final model. Dashed: the ideal width
  $\sqrt3/(\pi\snr)$. Every model has a floor at high SNR; the warped clock lowers it by a
  factor 0.6 and the cooldown by a further 0.44, and below SNR $\sim40$ the final model is
  within a factor 2--3 of the ideal.}
  \label{fig:width-snr}
\end{figure}

\subsection{What Each Change Bought}
\label{app:xs-ladder}

\Cref{fig:ladder} summarises the ladder. Detection differences are paired, source by source,
with a two-sided sign test.

\begin{figure}[h]
  \centering
  \includegraphics[width=\textwidth]{figures/ladder.pdf}
  \caption{The XS ladder. (a) The loud-source floor: median $f_0$ width at SNR $\ge100$, with
  its ratio to the ideal width. (b) Found fraction per SNR band. Same 840 sources for all
  models.}
  \label{fig:ladder}
\end{figure}

\paragraph{Training longer at a constant learning rate helps a little.}
Doubling the uniform-clock run from 500k to 1M steps narrowed the posteriors of the ten eval
windows by a median factor 0.78 (narrower for 27 of 29 sources) and localised two more of the
13 sources at SNR 15--40, but the loud floor only moved from 92 to 76 times the Fisher width.

\paragraph{The warped clock narrows every posterior, at a small cost in detection.}
Against the uniform clock at the same 1M steps, the warp narrowed the loud floor from 0.147 to
0.088 bins, by a median factor of 0.61 at SNR $\ge100$, 0.62 at 40--100 and 0.70 at
15--40, and the posterior was narrower for 97\% of the loud sources. It also removed a bias:
the uniform-clock model placed loud sources low by about 0.056 bins (mean rank 0.62, 8 standard
errors from 1/2). The price was detection: 27 sources were found only by the uniform clock and
6 only by the warp ($p = 3\times10^{-4}$), mostly near SNR 17, the expected effect of moving
training samples from the first half of the path (50\% of them) to its end (21\% left).

\paragraph{Width buys detection, not the floor.}
At a constant learning rate, the 768-wide network had the same floor as the 512 (0.087 against
0.088 bins, flat in SNR), with 2.25 times the parameters and a 1.3\% lower loss. It found 16
sources the 512 missed against 5 the other way ($p=0.027$), recovering half of the warp's
detection cost.

\paragraph{The cooldown halves the floor and removes the bias.}
Cooling the 768 network down over 200k more steps lowered the loss from 0.4499 to 0.4276; at
the constant-rate slope of its last 100k steps the same gain would have taken about a million
more steps. The loud floor fell from 0.087 to 0.037 bins, a median per-source ratio of 0.440
(95\% interval 0.427--0.453), at every loud SNR: 0.038, 0.039, 0.035 and 0.036 bins at SNR
100--200, 200--400, 400--800 and $\ge800$. The cooled model found 21 sources the 1M model did
not, and missed one ($p=10^{-5}$). Against the uniform clock it now finds more sources in every
band (17 against 7). The 512 network gained exactly as much from the same cooldown (factor
0.443), so the gain belongs to the schedule, not to the width. Once both are cooled, the floor
does not depend on width (the 512 is 1.03 times the 768, 95\% interval 1.01--1.05), but the 768
is 11--15\% sharper below SNR 100 and finds 18 sources the 512 misses against 7 ($p=0.043$).

\paragraph{What the floor is not.}
Cheap evaluation-only tests ruled out the obvious numerical causes (\cref{tab:numerics}): the
integrator (32, 64 and 128 RK4 steps agree), bfloat16 at evaluation (float32 on a CPU gives a
width ratio of 0.991), the rounding of the input coordinate at constant learning rate
(\cref{app:geometry-remarks}), and A100 numerics (the A100 and CPU tables agree). Width is ruled
out above. The floor that remains is flat in SNR at 25 times the ideal width. By
\cref{eq:sensitivity} it costs about $2\times10^{-3}$ of the logged loss under the warp, four
times the epoch-to-epoch scatter: visible to the objective, unlike the floor under the uniform
clock, but small next to the 0.022 the cooldown gained. Bfloat16 arithmetic inside the network
during training, and the optimisation itself, are the candidates left untested.

\subsection{Bias}

At a constant learning rate, each run carried its own smooth $f_0$ bias as a function of where
the source sits in the window (\cref{fig:bias}): up to $-0.146$ bins for the uniform clock and
$+0.081$ bins for the 768-wide warp, comparable to their widths, with a different shape in each
run. The cooldown removes it, to at most 0.008 bins across the window, and the mean
offset of the loud sources goes from $+0.029\pm0.005$ to $+0.007\pm0.005$ bins. A profile that
differs from run to run and disappears when the learning rate anneals is the signature of
where a noisy constant-rate iterate happened to stop, rather than of the physics or the data.

\begin{figure}[h]
  \centering
  \includegraphics[width=0.55\textwidth]{figures/bias.pdf}
  \caption{Median $f_0$ offset of the found sources at SNR $\ge40$ in eight slices of their
  position in the window, from its low edge ($-1$) to its high edge ($+1$); bands show one
  standard error of the median. At a constant learning rate each run has its own bias
  profile; after the cooldown none is left.}
  \label{fig:bias}
\end{figure}

\subsection{Calibration}

\Cref{fig:calibration} shows the calibration of the $f_0$ marginal over all 840 sources. The
uniform-clock model's ranks lean to one side, its bias. The final model's are centred but
over-represented in the middle: the truth falls near the centre of the posterior more often
than it should, the signature of posteriors that are wider than necessary, consistent with
in68 = 0.93 at SNR $\ge100$. No model is over-confident in any band. Calibration of the other
parameters will be measured against MCMC (\cref{app:mcmc}).

\begin{figure}[h]
  \centering
  \includegraphics[width=0.55\textwidth]{figures/calibration.pdf}
  \caption{Calibration of the $f_0$ marginal, 840 sources. (a) Empirical CDF of the truth's
  rank among the draws minus the uniform CDF, with a pointwise 95\% band. A calibrated
  posterior stays in the band; the uniform-clock model is biased, and the final model's
  S-shape means its posteriors are too wide. (b) Coverage of the central 68\% (circles) and
  95\% (squares) intervals by SNR band; dashed lines are the nominal levels.}
  \label{fig:calibration}
\end{figure}

\subsection{Single-Source Posteriors}

\Cref{fig:corner-median,fig:corner-loudest} show all eight parameters of two sources,
the median and the loudest of the 40 sources in the eval windows, against the Fisher forecast.
The flow reproduces the exact symmetries: four modes in $\psi$ at quarter-turn spacing, two in
$\phi_0$, and the diagonal bands they form together (\cref{eq:psi-phi-symmetry}). For the
loud source it follows the curved amplitude--inclination degeneracy towards face-on, where
$\psi$ and $\phi_0$ merge into one combination. In both cases the flow is wider than the Fisher
forecast in $f_0$ (the floor above) and along the amplitude--inclination direction. How much
of the latter is the true posterior, which a Gaussian forecast cannot follow, and how much is
excess width, is what the comparison with MCMC will measure (\cref{app:mcmc}). The chirp mass
at these frequencies is not measured, and the flow returns its prior.

\begin{figure}[p]
  \centering
  \includegraphics[width=0.92\textwidth]{figures/corner_median.pdf}
  \caption{Posterior of the median-SNR eval source (SNR 62.4, window at 0.888\,mHz), final
  model, 1024 draws relabelled onto the source (blue), against the Fisher forecast (grey) and
  the truth (black). Contours enclose 39\% and 86\% of the mass. The flow reproduces the four
  $(\psi,\phi_0)$ copies of \cref{eq:psi-phi-symmetry}. 18\% of the draws fall outside the
  plotted $f_0$ range: they are slots the relabelling assigns to this source although they
  belong to an unlocalised SNR-27 source in the same window (\cref{app:evaluation-matching}).}
  \label{fig:corner-median}
\end{figure}

\begin{figure}[p]
  \centering
  \includegraphics[width=0.92\textwidth]{figures/corner_loudest.pdf}
  \caption{As \cref{fig:corner-median}, for the loudest eval source (SNR 1366, window at
  0.101\,mHz). The $f_0$ posterior is far wider than the forecast (the floor); the amplitude
  and inclination follow the degeneracy towards face-on, where the polarisation angle and the
  phase become degenerate.}
  \label{fig:corner-loudest}
\end{figure}

\subsection{Training Dynamics}
\label{app:xs-dynamics}

\Cref{fig:training} shows the training curves. With the original long auxiliary phase, the
flow loss crept down at $2$--$3\times10^{-4}$ per epoch for 170 epochs and then fell by 0.056
within 35 epochs, as the auxiliary weight crossed about 0.2. With $f_{\rm aux}=0.1$ and the
uniform clock (a run stopped after 60 epochs to save compute), the crossover at epoch 36
brought no such drop, so the auxiliary weight alone does not trigger it. Under the warp the
drop came early, near epoch 60, consistent with the warp putting more weight on the fine
structure that the drop seems to learn; what triggers it remains open. After it, the loss
follows a slow power law: each doubling of the steps lowers it by 0.014 to 0.022, and a fit of
$L_\infty + A\,n^{-\alpha}$ to steps 300k--681k predicted 0.4070 at 1M steps, against the 0.4061
reached. The cooldown is worth much more than further constant-rate training:
$-0.022$ in 200k steps for the 768 network and $-0.020$ for the 512. The losses of the
uniform and warped clocks weight the path differently and are not comparable with each other.
Once the auxiliary weight reaches zero, the auxiliary heads stop training and their losses
drift up by orders of magnitude (the reconstruction loss reaches 190 and 520 at 1M steps, at
widths 512 and 768); they play no role in sampling.

\begin{figure}[h]
  \centering
  \includegraphics[width=\textwidth]{figures/training.pdf}
  \caption{Training curves: median flow loss per epoch of 1000 steps. (a) XS, uniform clock,
  500k steps then continued to 1M. (b) XS, warped clock, widths 512 and 768, each continued
  with a 200k-step cooldown (shaded). (c) B, warped clock, 512 wide: the run in progress
  (\cref{app:b}), with its planned cooldown. Losses under different clocks and problems are
  not comparable with each other.}
  \label{fig:training}
\end{figure}

\subsection{Numerical Checks}

\begin{table}[h]
\caption{The warped-clock model (512, 1M steps) under different integrators and precisions.
The float32 rows score the first 200 sources (50 windows) with the same keys; ``ratio'' is
the median per-source width ratio to the default bfloat16 evaluation on the same sources.
Autotuning on the A100 at level 3 gives over-confident posteriors (\cref{app:training-numerics}).}
\label{tab:numerics}
\centering
\small
\setlength{\tabcolsep}{4pt}
% generated by paper/scripts/tables.py from the scorecards; do not edit by hand
\begin{tabular}{lrccccccccc}
\toprule
& & \multicolumn{4}{c}{found fraction, $\rho$:} & \multicolumn{2}{c}{$f_0$ width, $\rho\geq100$} & \multicolumn{2}{c}{$\rho\geq100$} \\
\cmidrule(lr){3-6} \cmidrule(lr){7-8} \cmidrule(lr){9-10}
sampler & $n$ & $<15$ & 15--40 & 40--100 & $\geq100$ & bins & ratio & in68 & in95 \\
\midrule
bf16, A100, 32 RK4 steps (default) & 840 & 0.45 & 0.81 & 0.92 & 0.99 & 0.088 & 1.000 & 0.87 & 0.99 \\
bf16, A100, 64 steps & 840 & 0.45 & 0.81 & 0.92 & 0.99 & 0.089 & 1.000 & 0.86 & 0.99 \\
bf16, A100, 128 steps & 840 & 0.45 & 0.80 & 0.92 & 0.99 & 0.089 & 0.998 & 0.87 & 0.99 \\
fp32, CPU, 32 steps & 200 & 0.39 & 0.78 & 0.94 & 1.00 & 0.088 & 0.991 & 0.85 & 1.00 \\
fp32, A100, autotune level 0 & 200 & 0.39 & 0.78 & 0.94 & 1.00 & 0.088 & 0.991 & 0.85 & 1.00 \\
fp32, A100, autotune level 3 & 200 & 0.45 & 0.83 & 0.97 & 1.00 & 0.036 & 0.411 & 0.45 & 0.88 \\
\bottomrule
\end{tabular}

\end{table}

\subsection{Summary: the Recipe}

What the XS ladder settled, and what B uses: the warped clock with $p=3$, auxiliary heads for
the first 10\% of training, and a linear learning-rate cooldown over the last 20\%. Width 768
helps detection at 1.5 times the cost of 512 at equal steps. The remaining $f_0$ floor sits at
25 times the ideal width for loud sources, with posteriors that are conservative rather than
over-confident.

% =======================================================================================
\section{The Whole Band: B}
\label{app:b}
% =======================================================================================
% STATUS (2026-10-08): B1 is running on TREX (jobs 13882121 -> 13882126, 4 chained 24 h
% links + 1 spare). Epoch 91 of 200 at 07:45 UTC on 8 Oct, flow loss 0.526, 1536 s per
% epoch, no NaN; the cooldown starts at epoch 161, earliest finish ~08:30 CEST on 10 Oct.
% Everything marked [TBD] is filled from its scorecard and eval:
%   sbatch slurm/lisa-scorecard.sbatch B-late --n_random 50     # 60 windows, 240 sources, ~1 h
%   sbatch slurm/lisa-eval.sbatch B-late                         # corner pages, chunked draws
% then: copy scorecard.npz to paper/data/scorecards/B-late.npz, rerun paper/scripts.

\subsection{Setup}

B covers the whole band, 0.1--12\,mHz, in windows of 4096 bins ($R=1024$, an interior of 3072
bins) with 2048 data tokens, 64 times as many as XS. Everything else is the XS recipe
(\cref{app:xs}): four sources per window, the same priors, the 512-wide network (75.6M
parameters, unchanged), the warped clock with $p=3$, auxiliary heads for the first 10\% of the
steps, and a linear cooldown over the last 20\%. A step costs 1.54\,s on one A100 (41 times
XS) with a peak of 18.6\,GiB, so B can afford a fifth of the steps of the XS runs:

\begin{table}[h]
\caption{Cost of a B run on one A100, with the last 20\% of the steps cooling down, in 24-hour
jobs. At a fixed budget the 512-wide network gets 1.64 times the steps of the 768.}
\label{tab:b-cost}
\centering
\small
\begin{tabular}{lll}
\toprule
steps & width 512 & width 768 \\
\midrule
100k & 44\,h, 2 jobs & 71\,h, 3 jobs \\
150k & 65\,h, 3 jobs & 107\,h, 5 jobs \\
200k & 87\,h, 4 jobs & 143\,h, 6 jobs \\
\bottomrule
\end{tabular}
\end{table}

We train the 512-wide network for 200k steps, the last 40k cooling down. On XS, width 768
bought detection at equal steps (\cref{app:xs-ladder}), but B is far short of the XS step
counts, and each doubling of the steps was worth about 20\% in width on XS; at this budget,
steps are the better use of compute. \Cref{fig:training}c shows the run so far: the flow loss
fell from 0.753 to 0.526 in 91k steps and is still falling, with no instability. As on XS, the
auxiliary losses drift once their weight reaches zero (at 20k steps).

\subsection{Results}

\begin{table}[h]
\caption{B by SNR band, on 240 sources in 60 windows (\cref{app:evaluation}). \tbd}
\label{tab:b-band}
\centering
\small
\begin{tabular}{lrcccrccc}
\toprule
SNR $\snr$ & $n$ & found & width [bins] & $\times$ ideal & offset [bins] & in68 & in95 & rank \\
\midrule
$<15$ & \tbd & \tbd & \tbd & \tbd & \tbd & \tbd & \tbd & \tbd \\
15--40 & \tbd & \tbd & \tbd & \tbd & \tbd & \tbd & \tbd & \tbd \\
40--100 & \tbd & \tbd & \tbd & \tbd & \tbd & \tbd & \tbd & \tbd \\
$\geq 100$ & \tbd & \tbd & \tbd & \tbd & \tbd & \tbd & \tbd & \tbd \\
\midrule
all & \tbd & \tbd & & & & \tbd & \tbd & \tbd \\
\bottomrule
\end{tabular}
\end{table}

\begin{table}[h]
\caption{B by frequency, in thirds of the log band. Above a few mHz the chirp is resolved and
the chirp mass becomes measurable; the Doppler sideband widens and the sky position sharpens.
Chirp-mass and sky widths are for found sources at SNR $\ge40$. \tbd}
\label{tab:b-frequency}
\centering
\small
\begin{tabular}{lccccc}
\toprule
$f_0$ & found ($\snr\ge40$) & loud $f_0$ width [bins] & $\sigma_{\ln\Mc}$ & sky width [deg] & in95 \\
\midrule
0.1--0.49\,mHz & \tbd & \tbd & \tbd & \tbd & \tbd \\
0.49--2.4\,mHz & \tbd & \tbd & \tbd & \tbd & \tbd \\
2.4--12\,mHz & \tbd & \tbd & \tbd & \tbd & \tbd \\
\bottomrule
\end{tabular}
\end{table}

\begin{figure}[h]
  \centering
  \begin{minipage}{0.48\textwidth}\centering
    \placeholderfig[0.18\textheight]{$f_0$ width against SNR on B, with the XS final model and
    the ideal width (as \cref{fig:width-snr}).}
  \end{minipage}\hfill
  \begin{minipage}{0.48\textwidth}\centering
    \placeholderfig[0.18\textheight]{Corner plot of a loud high-frequency source, where the
    chirp mass is measured (as \cref{fig:corner-loudest}).}
  \end{minipage}
  \caption{B results. \tbd}
  \label{fig:b-results}
\end{figure}

The analysis that will fill these placeholders:
\begin{itemize}
  \item \textbf{Does the recipe carry over?} Detection, widths and calibration by SNR band
        (\cref{tab:b-band}), against the XS final model on the same SNR bands. \tbd
  \item \textbf{What the band adds.} At high frequency the chirp is resolved and the chirp
        mass measured (\cref{tab:b-frequency}); the Fisher forecast is then informative for
        $\Mc$ and is the reference for its width. \tbd
  \item \textbf{Does the wider window cost precision?} The B and XS models, compared on
        windows below 1.26\,mHz, where both are trained. \tbd
  \item \textbf{Resolution of the frequency coordinate.} One unit of $x_f$ spans 1536 bins on
        B, so its bfloat16 cell reaches 6 bins over half of the window
        (\cref{app:geometry-remarks}). The $f_0$ width against the source's position in the
        window, as in \cref{fig:bias}, tests whether this limits B. \tbd
\end{itemize}

% Predictions written before the result (2026-10-08), for the authors; not for the paper:
%  - the loss drops at the cooldown (epoch 161) as on XS, by ~0.02;
%  - if the recipe carries over: loud f0 floor near XS's 0.04 bins (25-30x ideal), SNR 15-40
%    detection >= 80%, in95 >= 0.95 in every band;
%  - if the bf16 cell of x_f matters on B: the loud width grows with |x_f| and sources with
%    |x_f| > 0.5 (half of them) are wider than a bin, so loud-source detection falls well
%    below XS's 99%. On XS the effect is at most 1.1-1.2x at a 0.094-bin cell; B's cell is
%    64x larger in bins. The fix would be fp32 Fourier features of the source coordinates
%    before the bf16 cast, which needs a new run.

% =======================================================================================
\section{Comparison with MCMC}
\label{app:mcmc}
% =======================================================================================
% STATUS (2026-10-10): filled. lisa_checks/mcmc_xs.py (jexplore, needs .venv313) ran on a laptop
% GPU, two seeds per window; lisa_checks/compare_mcmc.py makes the per-window summaries, which
% scripts/import_data.py mcmc copies into data/mcmc. Tables from scripts/tables.py, the figure
% from scripts/fig_mcmc.py. Research log F27.

\subsection{MCMC for Galactic Binaries}

The standard analysis of LISA's galactic binaries is a global fit
\citep{cornish2007tests,littenberg2020global,littenberg2023prototype,strub2024global,katz2024efficient,deng2025modular}:
a Markov chain samples all resolvable sources at once, together with the noise, moves
between models with different numbers of sources by reversible jumps
\citep{green1995reversible}, and uses parallel tempering
\citep{earl2005parallel,vousden2016dynamic} to cross between the modes of
\cref{app:priors-symmetries}. The frequency band is split into overlapping segments
that are updated in turn, which is what makes the problem tractable and is also the unit our
windows imitate. Each new stretch of data, or each change of the noise model, means sampling
again. An amortised posterior estimator costs one forward integration per window once
trained, which is what makes it attractive as a proposal distribution, as a fast first pass, or
for population studies that need many independent analyses \citep{srinivasan2025simulation}.
Neural density estimation for galactic binaries was explored by \citet{korsakova2024neural}.

\subsection{Sampler}

Our reference sampler is \texttt{jexplore}, a JAX parallel-tempered ensemble MCMC developed in
the LISA community, with affine-invariant stretch moves \citep{goodman2010ensemble,foreman2013emcee} and temperature swaps, run on the likelihood of \cref{eq:loglik} and the
priors of \cref{tab:prior}. For one source (8 dimensions), 32 walkers at four temperatures
(1, 1.6, 2.56, 4.1) for 20\,000 iterations, the second half kept, take about a minute on a GPU,
with acceptance rates of 0.24 (stretch) and 0.32 (swap) and $\hat R<1.0002$. On a published
source, this sampler on our likelihood reproduces the published posterior
(\cref{app:noise-validation}).

For the four-source windows below, the sampler runs on all 32 parameters jointly. Each source
is sampled as $(\ln f_0, \ln\Mc, \ln\Amp, \alpha, \sin\delta, \psi, \cos\iota, \phi_0)$ under the
training prior (\cref{tab:prior}), with the likelihood in double precision. The three angles are
sampled on $2\pi$ intervals centred on their true values. That is a full period, so the prior
is unchanged, and no narrow mode is cut by an edge. The chirp-mass prior has no closed form, so
we compute its density by quadrature over the primary mass. It matches $2\times10^6$ prior
draws ($\chi^2/\text{dof}=0.72$ over 60 bins).

Each iteration is a temperature swap, followed by either a differential-evolution or a stretch
move, with equal probability. There are 64 walkers at each of six temperatures from 1 to 10.
The walkers start from the window's Fisher draws, with chirp masses drawn from the prior. We
discard 100\,000 iterations and keep 300\,000, thinned by 100, and run each window twice with
independent seeds.

\subsection{Flow against MCMC}

\paragraph{Protocol.}
Two XS eval windows, chosen to probe the two limitations of \cref{app:xs}:
\begin{itemize}
  \item the loudest (window SNR 1848, sources at SNR 14, 33, 1245 and 1366), where the $f_0$
        floor and the amplitude--inclination width matter;
  \item the crowded window of \cref{fig:corner-median} (sources at SNR 26, 27, 62 and 567),
        where the flow misses an SNR-27 source.
\end{itemize}
For each window, the MCMC uses the same noise realisation as the flow. We relabel its samples
as the flow's (\cref{app:evaluation-matching}), and fold $\psi$ and $\phi_0$ of both sample
sets by the exact symmetries of \cref{eq:psi-phi-symmetry} before comparing. For all eight
parameters we then compare:
\begin{itemize}
  \item the marginal widths ($1.4826\,\MAD$);
  \item the 1-D Wasserstein distances, in units of the MCMC width;
  \item the coverage of the MCMC medians by the flow's intervals.
\end{itemize}

\paragraph{The MCMC is converged, and Fisher holds where the floor is measured.}
\Cref{tab:mcmc-checks} checks the sampler itself.
\begin{itemize}
  \item The two seeds agree within 7.5\% in every width and within $0.1\sigma$ in
        $W_1$. The exception is the SNR-14 source ($0.49\sigma$), whose sky posterior is broad
        and multimodal.
  \item At SNR $\ge567$, the Fisher forecast matches the MCMC's $f_0$ width to 2\%, and every
        other width within a factor 0.8--1.4. So the ``$\times$ ideal'' and ``$\times$
        Fisher'' of \cref{app:xs} compare the flow with the true posterior where the floor is
        quoted.
  \item Below SNR $\sim60$, the forecast can be off by a factor of two: along the
        amplitude--inclination direction near face-on, and on the sky of the SNR-14 source.
\end{itemize}

\begin{table}[h]
\caption{The MCMC itself, per source: its $f_0$ width, the Fisher forecast against it, and the
two seeds against each other. ``Other'' spans the six parameters besides $f_0$ and the
prior-dominated $\Mc$. The seed columns are maxima over the eight parameters: of
$|$width ratio $-1|$, and of $W_1$ in units of the MCMC width.}
\label{tab:mcmc-checks}
\centering
\small
% generated by paper/scripts/tables.py from data/mcmc; do not edit by hand
\begin{tabular}{lrccccc}
\toprule
& & MCMC $f_0$ & \multicolumn{2}{c}{width, Fisher / MCMC} & \multicolumn{2}{c}{seed 1 against seed 0} \\
\cmidrule(lr){4-5} \cmidrule(lr){6-7}
window & SNR & width [bins] & $f_0$ & other & max $|$width ratio $-1|$ & max $W_1/\sigma$ \\
\midrule
loudest & 1365.6 & 0.00045 & 0.98 & 0.81--1.39 & 0.031 & 0.04 \\
loudest & 1245.4 & 0.00056 & 0.98 & 0.79--1.04 & 0.075 & 0.10 \\
loudest & 32.9 & 0.016 & 1.03 & 0.93--1.42 & 0.042 & 0.05 \\
loudest & 14.4 & 0.046 & 0.93 & 0.40--0.99 & 0.057 & 0.49 \\
\midrule
crowded & 567.4 & 0.0013 & 0.98 & 0.92--1.05 & 0.020 & 0.02 \\
crowded & 62.4 & 0.0097 & 0.99 & 0.49--1.27 & 0.013 & 0.02 \\
crowded & 27.1 & 0.023 & 1.06 & 0.67--1.12 & 0.012 & 0.02 \\
crowded & 26.3 & 0.023 & 0.95 & 0.93--1.07 & 0.012 & 0.01 \\
\bottomrule
\end{tabular}

\end{table}

\paragraph{Flow against MCMC.}
\Cref{tab:mcmc} gives the width ratios, and \cref{fig:mcmc} the loudest source.
\begin{itemize}
  \item \textbf{Never biased.} The MCMC median lies inside the flow's central 95\% interval for
        all 64 source--parameter pairs, and inside its 68\% interval for 57 of them. The flow is
        conservative, as the coverage of \cref{app:xs} said.
  \item \textbf{Loud sources are too wide in every parameter.} $f_0$ is 38--84 times the MCMC
        width, the floor of \cref{app:xs}. The sky, amplitude, inclination and the two
        phases are 4--8 times too wide. The MCMC follows the same curved
        amplitude--inclination degeneracy as the flow (\cref{fig:mcmc}), but is four times
        narrower along it. So the flow's excess width there is not a non-Gaussianity that
        the Fisher forecast misses, and the precision limit is not specific to $f_0$, though
        $f_0$ is where it costs most.
  \item \textbf{Faint sources, once found, are near-optimal.} At SNR 26--33 the flow is within
        1.0--2.1 times the MCMC in every parameter but $f_0$ (2.3--4.8).
  \item \textbf{One faint source per window is missed, by the flow, not by the data.} The MCMC
        localises the SNR-27 source of the crowded window to 0.023 bins, as the Fisher
        forecast predicts, while 87\% of the flow's draws for it are more than a bin away.
        The SNR-14 source of the loudest window is similar (58\%). When one slot does not
        lock on, a brighter neighbour loses about a fifth of its draws to the confusion (SNR
        33 and 62).
  \item $\psi$ and $\phi_0$ of the SNR-62 source are nearly face-on ($\cos\iota\approx0.95$).
        Only a combination of them is constrained, in both methods, so the ratio of their
        separate widths (0.5) means little; $W_1$ is $0.3\sigma$. The chirp mass is
        prior-dominated in both (ratios 0.95--1.06).
\end{itemize}

\begin{table}[h]
\caption{Flow against MCMC on the two windows: ratio of marginal widths (flow / MCMC) per
parameter, after relabelling and folding $\psi$ and $\phi_0$ by their exact symmetries. The
last column is the share of the flow's draws more than a bin from the MCMC median in $f_0$. In
every entry the MCMC median lies inside the flow's central 95\%.}
\label{tab:mcmc}
\centering
\small
% generated by paper/scripts/tables.py from data/mcmc; do not edit by hand
\begin{tabular}{lrccccccccr}
\toprule
& & \multicolumn{8}{c}{width, flow / MCMC} & off \\
\cmidrule(lr){3-10}
window & SNR & $f_0$ & $\Mc$ & $\Amp$ & $\alpha$ & $\sin\delta$ & $\psi$ & $\cos\iota$ & $\phi_0$ & $>1$ bin \\
\midrule
loudest & 1365.6 & 84 & 0.95 & 4.21 & 5.44 & 5.72 & 6.62 & 4.29 & 6.62 & 0\% \\
loudest & 1245.4 & 66 & 1.06 & 3.62 & 6.37 & 6.41 & 8.42 & 3.86 & 8.39 & 0\% \\
loudest & 32.9 & 4.78 & 0.95 & 1.36 & 1.76 & 2.06 & 1.06 & 2.11 & 1.03 & 20\% \\
loudest & 14.4 & 21 & 0.97 & 5.14 & 4.09 & 5.41 & 2.04 & 8.38 & 9.51 & 58\% \\
\midrule
crowded & 567.4 & 38 & 0.99 & 4.81 & 5.75 & 4.47 & 5.00 & 4.11 & 4.76 & 0\% \\
crowded & 62.4 & 6.24 & 1.02 & 2.58 & 6.76 & 3.77 & 0.47 & 1.71 & 0.53 & 18\% \\
crowded & 27.1 & 231 & 1.01 & 4.45 & 54 & 34 & 2.77 & 3.66 & 5.59 & 87\% \\
crowded & 26.3 & 2.30 & 1.04 & 1.03 & 1.12 & 1.20 & 1.17 & 1.08 & 1.15 & 0\% \\
\bottomrule
\end{tabular}

\end{table}

\begin{figure}[h]
  \centering
  \includegraphics[width=0.92\textwidth]{figures/mcmc.pdf}
  \caption{The loudest eval source (SNR 1366, window at 0.1009\,mHz): the flow (blue), a joint
  MCMC of the window's four sources (black), the Fisher forecast (grey, dashed) and the truth
  (red). Contours enclose 39\% and 86\% of the mass. The MCMC draws are moved at random to
  the four exact $(\psi,\phi_0)$ copies, to show them on the flow's footing. MCMC and Fisher
  agree. The flow covers both, but is 4--8 times wider in every parameter and 84 times wider in
  $f_0$.}
  \label{fig:mcmc}
\end{figure}

\paragraph{Cost.}
Once trained, the flow draws 256 posterior samples of an XS window in about a second on an
A100, and of a B window in about 42 seconds. The MCMC of one XS window takes 14 minutes per
seed on a laptop GPU: 400\,000 iterations of 384 chains with a double-precision likelihood.
Training costs about 20 A100-hours on XS and 86 on B, paid once.


\end{document}

% This document was modified from the file originally made available by
% Pat Langley and Andrea Danyluk for ICML-2K. This version was created
% by Iain Murray in 2018, and modified by Alexandre Bouchard in
% 2019 and 2021 and by Csaba Szepesvari, Gang Niu and Sivan Sabato in 2022.
% Modified again in 2023 and 2024 by Sivan Sabato and Jonathan Scarlett.
% Previous contributors include Dan Roy, Lise Getoor and Tobias
% Scheffer, which was slightly modified from the 2010 version by
% Thorsten Joachims & Johannes Fuernkranz, slightly modified from the
% 2009 version by Kiri Wagstaff and Sam Roweis's 2008 version, which is
% slightly modified from Prasad Tadepalli's 2007 version which is a
% lightly changed version of the previous year's version by Andrew
% Moore, which was in turn edited from those of Kristian Kersting and
% Codrina Lauth. Alex Smola contributed to the algorithmic style files.
